%%%%%%%%%%%%%%%%%%%%%%%%%%%%%%%%%%%%%%%%%%%%%%%%%%%%%%%%%%%%%%%%%%%%%%%%%%
%                                                                        %
%                         Chapter 9 - Section 1                          %
%                                                                        %
%%%%%%%%%%%%%%%%%%%%%%%%%%%%%%%%%%%%%%%%%%%%%%%%%%%%%%%%%%%%%%%%%%%%%%%%%%

\chapter{Functions of Several Variables}

Functions that depend on several input variables first appeared in the
$S$-$I$-$R$ model at the beginning of the course.  Usually, the number
of variables has not been an issue for us.  For instance, when we
introduced the derivative in chapter 3, we used partial derivatives to
treat functions of several variables in a parallel fashion.
%
\mar{The problem of visualizing a function of several variables}	
%
However, when there are questions of visualization and geometric
understanding, the number of variables {\em does\/} matter.  Every
variable adds a dimension to the problem---one way or another.  For
example, if a function has two input variables instead of one, we will
see that its graph is a surface rather than a curve.
		
This chapter deals with the geometry of functions of two or more
variables.  We start with graphs and level sets.  These are the basic
tools for visualization.  Then we turn to microscopic views, and see
what form the microscope equation takes.  Finally, we consider
optimization problems using both direct visual methods and dynamical
systems.
	
	
	
	
\section{Graphs and Level Sets}

\special{psfile=../ps/911/WEATHER1.PS}

\begin{minipage}[t]{3.5in}
The graph\index{graph} at the right comes from a model that describes
how the average daily temperature\index{temperature function} at one
place varies over the course of a year.  It shows the temperature $A$
in $^\circ$F, and the time $t$ in months from January.  As we would
expect, the temperature is a periodic function (which we can write as
$A(t)$), and its period is 12 months.  Furthermore, \linebreak
\end{minipage}  

\noindent
the lowest temperature occurs in February (when $t \approx$ 1 or 13)
and the highest in July (when $t \approx$ 7 or~19).  This is about
what we would expect.

However, all these temperature fluctuations disappear 
%
\mar{Underground temperatures\\ fluctuate less}	
%
a few feet underground.  Below a depth of 6 or 8 feet, the temperature
of the soil remains about 55$^\circ$F year-round!  Between ground
level and that depth, the temperature still fluctuates, but the range
from low to high decreases with the depth.  Here is what happens at
some specific depths.

 

\special{psfile=../ps/911/WEATHER2.PS}
\vspace{166pt}
\label{timegraph}

\noindent
Notice how the time at which the temperature peaks gets later and
later as we go farther and farther underground.  For example, at $d =
2$ feet the highest temperature occurs in September ($t \approx 9$),
%
\mar{The phase shifts\\ with the depth}
%
not July.  It literally takes time for the heat to sink in.  In
chapter 7 we called this a phase shift.\index{phase!shift} The lowest
temperature shifts in just the same way.  At a depth of 2 feet, it is
colder in March than in January.

Thus $A$ is really a function of {\em two\/}
variables,\index{function!of two variables} the depth $d$ as well as
the time $t$.  To reflect this addition, let's change our notation for
the function to $A(t, d)$.  In the figure above, $d$ plays the role of
a parameter: it has a fixed value for each graph.  We can reverse
these roles and make $t$ the parameter.  This is done in the figure on
the top of the next page.  It shows us how the temperature
%
\mar{Graphing temperature as a function of depth}	
%
varies with the depth at fixed times of year. Notice that, in April
and October, the extreme temperature is not found on the surface.  In
October, for example, the soil is warmest at a depth of about 9
inches.
	
The lower figure on the page is a single graph that combines all the
information in these two sets of graphs.  Each point on the bottom of
the box of the box corresponds to a particular depth and a particular time.
The height of the surface above that point tells us the temperature at
that depth and time.\index{graph!reading a}\index{graph!surface}
For example, suppose you want to find the temperature 4 feet
\linebreak

\mbox{}

\special{psfile=../ps/911/WEATHER3.PS}
\vspace{150pt}
\label{depthgraph}
	
\noindent
below surface at the beginning of July. 
%
\mar{Reading a\\ surface graph}
%
Working from the bottom front corner of the box, move 4 feet to the
right and then 6 months toward the back.  This is the point $(d, s) =
(4, 6)$.  The height of the graph above this point is the temperature
$A$ that we want.
    	
\special{psfile=../ps/911/WEATHER5.PS}
	
\newpage

There is a definite connection between this surface and the two
collections of curves.
%
\mar{Grid lines are slices of the surface that show how the function depends on each variable separately}	
%
Imagine that the box containing the surface graph is a loaf of bread.
If you slice the loaf parallel to the left or right side, this slice
is taken at a fixed depth.  The cut face of the slice will look like
one of the graphs on page~\pageref{timegraph}.  These show how the
temperature depends on the time at fixed depths.  If you slice the
loaf the other way---parallel to the front or back face---then the
time is fixed.  The cut face will look like one of the graphs on
page~\pageref{depthgraph}.  They show how the temperature depends on
the depth at fixed times.  The grid lines\index{grid line} on the
surface are precisely these ``slice''\index{slice}\index{graph!slice
  of a} marks.


\special{psfile=../ps/911/WEATHER6.PS}
\hfill
\begin{minipage}[t]{2.7in}
\quad \ Here is the same graph seen from a different viewpoint.  Now
time is measured from the left, while depth is measured from the back.
The temperature is still the height, though.  One advantage of this
view is that it shows more clearly how the peak temperature is
phase-shifted with the depth.

\quad \ We now have two ways to visualize how the average daily
temperature depends on the time of year and the depth below ground.
One is the surface graph itself, and the other is a collection of
curves that are slices of the surface.  The surface gives us an
\linebreak%
\end{minipage}

\vspace{-9pt}

\noindent
overall view, but it is not so easy to read the surface graph to
determine the temperature at a specific time and depth.
%
\mar{Comparing the surface to its slices}
%
Check this yourself: what is the temperature 2 feet underground at the
beginning of April?  The slices are much more helpful here.  You
should be able to read from either collection of slices that $A
\approx 44^\circ$F.

\subsection{Examples of Graphs}

\special{psfile=../ps/911/MINMAX1.PS}

The purpose of this section is to get some experience constructing and
interpreting surface graphs.  To work in a context, look first at the
functions $y = x^2$ and $y = -x^2$.  They provide us with standard
examples of a minimum and a maximum when there is just one input
variable.  Let's consider now the corresponding examples for two input
variables.  Besides an ordinary maximum and an ordinary minimum, we
will find a {\em third\/} type---called a minimax---that is completely
new.  It arises because a function can have a minimum with respect to
one of its input variables and a maximum with respect to the other.
	
\subsubsection{A minimum: $z = x^2 + y^2$}
\index{minimum}

\special{psfile=../ps/911/MIN2.PS}

\begin{minipage}[t]{3.7in}
\quad \ At the origin $(x, y) =(0, 0)$, $z = 0$.  At any other point,
either $x$ or $y$ is non-zero. Its square is positive, so $z > 0$.
Consequently, $z$ has a minimum at the origin.  The graph of this
function is a parabolic {\bf bowl} whose lowest point sits on the
origin.  As always, the grid lines are slices, made by fixing the
value of $x$ or $y$.  For example, if $y = c$, then the slice is $z =
x^2 + c^2$.  This is an ordinary parabolic curve.
\end{minipage}
		
		
\subsubsection{A maximum: $z = -x^2 - y^2$}
\index{maximum}

\special{psfile=../ps/911/MAX2.PS}

\begin{minipage}[t]{3.7in}
\quad \ For any $x$ and $y$, the value of $z$ in this example is the
opposite of its value in the previous one.  Thus, $z$ is everywhere
negative, except at the origin, where its value is 0.  Thus $z$ has a
maximum at the origin.  Its graph is an upside-down bowl, or {\bf
  peak}, whose highest point reaches up and touches the origin.  Grid
lines are the curves $z = -x^2 - c^2$ and $z = -c^2 - y^2$.  These are
parabolic curves that open downward.
\end{minipage}

\subsubsection{A minimax: $z = x^2 - y^2$}
\index{minimax}

\special{psfile=../ps/911/SADDLE1.PS}

\label{saddle}
\begin{minipage}[t]{3.7in}
\quad \ Suppose we fix $y$ at $y = 0$.  This slice has the equation $z
= x^2$, so it is an ordinary parabola (that opens upward).  Thus, as
far as the input $x$ is concerned, $z$ has a {\em minimum\/} at the
origin.  Suppose, instead, that we fix $x$ at $x = 0$. Then we get a
slice whose equation is $z = -y^2$.  It is also an ordinary parabola,
but this one opens downward.  As far as $y$ is concerned, $z$ has a
{\em maximum\/} at the origin.  It is \linebreak
\end{minipage}
\vspace{-9pt}

\noindent
clear from the graph how upward-opening slices in the $x$-direction
fit together with downward-opening slices in the $y$-direction.
Because of the shape of the surface, a minimax is commonly called a
{\bf saddle}, or a {\bf saddle point}.\index{saddle}\index{saddle
  point}

\newpage

Here are two slices\index{slice}\index{graph!slice of a} of $z = x^2
- y^2$ shown in more detail.  Points in the box have three
coordinates: $(x, y, z)$.  If we set $y = 0$ we are selecting the points
\linebreak

\special{psfile=../ps/911/SADDLE2.PS}

\special{psfile=../ps/911/SADDLE3.PS}

\special{psfile=../ps/911/SADDLE4.PS}
\vspace{210pt}

\noindent
\begin{minipage}{3.6in}
of the form $(x, 0, z)$.  These make up the $x, z$-plane.  On this
plane the equation $z = x^2 - y^2$ becomes simply $z = x^2$.  The
graph of {\em this\/} equation is a curve in the $x,
z$-plane---specifically, the parabola shown.  The situation is similar
if $y$ is given some\linebreak 
\end{minipage}

\vspace{-9pt}

\noindent
other fixed value. For example, $y = -4$ specifies the points $(x, -4,
z)$.
%
\mar{\vspace{-36pt} The points where $y = c$ form a vertical plane parallel to the $x, z$-plane}
%
These describe the plane that forms the front face of the box.  The
equation $z = x^2 - y^2$ becomes $z = x^2 - 16$.  The curve tracing
out the intersection of the saddle with the front of the box is
precisely the graph of $z = x^2 - 16$.

If $x = 0$ we get the points $(0, y, z)$ that make up the $y,
z$-plane.  On this plane the equation simplifies to $z = -y^2$, and
its graph is the parabolic curve shown.
%
\mar{The points where $x = c$ form a vertical plane parallel to the
  $y, z$-plane}
%
Giving $x$ a different fixed value leads to similar results.  A good
example is $x = -4$.  The points $(-4, y, z)$ lie on the plane that
forms the left side of the box.  The equation becomes $z = 16 - y^2$
there, and this is the parabolic curve marking the intersection of the
saddle with the left side of the box.

\aside{As you can see, it is valuable for you to be able to generate
  surface graphs yourself.  There are now a number of computer
  utilities which will do the job.  Some can even rotate the surface
  while you watch, or give you a stereo view.  However, even without
  one of these powerful utilities, you should try to generate the
  slicing curves that make up the grid lines of the surface.}

\subsubsection{A cubic: $z = x^3 - 4x - y^2$}
\label{cubic}

Slices of this graph are downward-opening parabolas (when $x = c$) and
are cubic curves that have the same shape (when $y = c$).  Notice that
each cubic curve has a maximum and a minimum, and each parabola has
\linebreak

\special{psfile=../ps/911/CUBIC1.PS}

\special{psfile=../ps/911/CUBIC2.PS}

\vspace{110pt}

\noindent
a maximum.  
%
\mar{The surface\\ has a saddle}
%
The surface graph itself has a {\em peak\/} where the cubics have
their maximum, but it has a {\em saddle\/}\index{saddle} where the
cubics have a minimum.  Do you see why?  The saddle point is a minimax
for $z = x^3 -4x -y^2$: $z$ has a minimum there {\em as a function of
  $x$ alone} but a maximum {\em as a function of $y$ alone}.

\vspace{10pt}

\special{psfile=../ps/911/CUBIC3.PS}
\special{psfile=../ps/911/CUBIC4.PS}
\special{psfile=../ps/911/CUBIC5.PS}
\vspace{230pt}

The small figures on the right show the same surface as the large
%
\mar{See the graph from different viewpoints}	
%
figure; they just show it from different viewpoints.  As a practical
matter, you should look at these surfaces the way you would look at
sculpture: ``walk around them'' by generating diverse views.
	
\subsubsection{Energy of the pendulum: $E = 1 - \cos{\theta} + \tfrac{1}{2} v^2$}

\special{psfile=../ps/911/PENDEN1.PS}
\vspace{200pt}

\special{psfile=../ps/911/PENDEN2.PS}
\hfill
\begin{minipage}[t]{3.5in}
\quad \ This function first came up in chapter 7, where it was used to
demonstrate that a dynamical system describing the motion of a
frictionless pendulum had periodic
solutions.\index{pendulum}\index{energy} It was used again in chapter
8 to clarify the phase portrait of that dynamical system.  The
function $E$ varies periodically with $\theta$, and you can see this
in the graph.  The minimum\index{minimum}\index{energy minimum} at the
origin is repeated at $(\theta, v) = (2\pi, 0)$, and so on.  The graph
also has a saddle at the point $(\theta, v) = (-\pi, 0)$.  This too
repeats with period $2\pi$ in the $\theta$ direction.

\quad \ The figure at the left is the same surface with part cut away
by a slice of the form $v = c$.  These slices are sine curves: $E = 1
- \cos{\theta} + \tfrac{1}{2} c^2$.  Slices of the form $\theta = c$
are upward-opening parabolas.  From this viewpoint, the saddle points
show up clearly.
\end{minipage}

\aside{One way to describe what happens to a real pendulum---that is,
  one governed by frictional forces as well as gravity---is to say
  that its energy ``runs down'' over time.  Now, at any moment the
  pendulum's energy is a point on this graph. As the energy runs down,
  that point must work its way down the graph.  Ultimately, it must
  reach the bottom of the graph---the minimum energy point at the
  origin $(\theta, v) = (0, 0)$.  This is the stable equilibrium
  point.  The pendulum hangs straight down ($\theta = 0$) and is
  motionless ($v = 0$).  The graph gives us an abstract---but still
  vivid and concrete---way of thinking of the dissipation of
  energy.\index{energy!dissipation of}}
	

\subsection{From Graphs to Levels}

\special{psfile=../ps/912/DENSITY2.PS}
\special{psfile=../ps/912/DENSITY1.PS}

\begin{minipage}[t]{3.5in}
There is still another way to picture a function of two variables.  To
see how it works we can start with an ordinary graph.  On the right is
the graph of
\[
z = f(x, y) = x^3 - 4x - y^2,
\]
the cubic function we considered on page~\pageref{cubic}.  This graph
looks different, though.  The difference is that points are shaded
according to their height.  Points at the bottom are lightest, points
at the top are darkest.

\quad \ Notice that the flat $x, y$-plane is shaded exactly like the
graph above it.  For instance, the dark spot centered at the point
$(x, y) = (-1, 0)$ is directly under the peak on the graph.  The other
\linebreak
\end{minipage}

\vspace{-10pt}

\noindent
dark patch, near the right edge of the plane, is under the highest
visible part of the surface.  Consequently, the shading on the $x,
y$-plane gives us the same information as the graph.  In other words,
{\em the intensity of shading at $(x, y)$ is proportional to the value
  of the function $f(x, y)$}.

\label{standard min max-begin}
The figure in the $x, y$-plane is called a {\bf density
  plot}.\index{plot!density}
%
\mar{Density plots}	
%
Think of the intensity of shading as the {\em
  density\/}\label{density} of ink on the page.  Here are density
plots of the standard minimum,\index{minimum} maximum,\index{maximum}
and minimax.\index{minimax} Compare these with the \linebreak
	
\special{psfile=../ps/912/DENSITY3.PS}
\vspace{140pt}

\noindent
\mbox{
	\small
	\mbox{} \hspace{25pt}
	$z = x^2 + y^2$
	\hspace{70pt}
	$z = -x^2 - y^2$
	\hspace{80pt}
	$z = x^2 - y^2$
}

\vspace{6pt}

\noindent
graphs on page~\pageref{saddle}.  The third density plot is the most
interesting.  From the center of the $x, y$-plane, the shading
increases to the right and left.  Therefore, $z$ has a minimum in the
horizontal direction.  However, the shading {\em decreases\/} above
and below the center.  Therefore, $z$ has a {\em maximum\/} in the
vertical direction.  Thus, you really can see there is a minimax at
the origin.

Try your hand at reading the density plot on the left below.  
%
\mar{A sample plot}	
%
You should see two maxima (directly above and below the origin), a
minimum (at the origin itself), and two saddles (to the right and the
left of the origin).  The function defining the plot is
\[
f(x, y) = (x^2 + (y-1)^2 - 3)(3 - x^2 - (y+1)^2).
\]

\special{psfile=../ps/912/DENSITY4.PS}
\special{psfile=../ps/912/density5.ps}
\vspace{210pt}

\noindent
Can you visualize what the graph looks like?  This density plot should
help you, and you can also construct slices by setting $x = c$ and $y
= c$.  The slices $x = 0$ and $y = 0$ are especially useful.  With
them you could determine the exact coordinates of the maxima and the
saddles.

\aside{These density plots show a ``checkerboard'' pattern because
  {\em Mathematica\/}\index{Mathematica} (the computer program that
  produces them) shades each little square according to the value of
  the function at the center of the square.  This pattern is an
  artefact; it is not inherent to a density plot.}

In a density plot, the shading varies smoothly with the value of the
function.  This is accurate, but it may be a bit difficult to read.
On the right you see a modified density plot.  There is still shading,
but there are now just a few distinct shades.
%
\mar{From densities \\to contours}
%
This makes a sharp boundary between one shade and the next.  The
boundary is called a {\bf contour},\index{contour} or a {\bf
  level}.\index{level} The figure itself is called a {\bf contour
  plot}.\index{plot!contour} The two maxima on the vertical line $x =
0$ stand out more clearly on the contour plot.  Also, the contour
lines around the two saddles help us see that the function has a
minimum in the vertical direction and a maximum in the horizontal
direction.
 

\special{psfile=../ps/913/CONTOUR1.PS}  

\noindent
\begin{minipage}[t]{3.4in}
\quad \ Once we have contour lines to separate one density level from
the next, we can even dispense with the shading.  The figure on the
right is just the contour plot from the opposite page, minus the
shading.  The contour lines, or level curves, now stand out clearly.
On each contour, the value of the function is constant.  This is also
called a {\bf contour plot}.

\quad \ There is some loss of information here, however.  For example,
we can't tell where the value of the function is large and where it is
small.  Nevertheless, the nested ovals on the vertical line $x = 0$
{\em do\/} tell us that there is either a maximum or a minimum at the
center of each nest.
\end{minipage}

\vspace{5pt}

For reference purposes, here are the contour plots for the standard
minimum,\index{minimum} maximum,\index{maximum} and
saddle.\index{saddle} In the first two cases, the contours are
concentric ovals.
%
\mar{Contours of the standard functions}
%
These look the same, so only one is illustrated.  The other two
pictures show a saddle.  In general, the contours around a saddle are
a family of hyperbolas.  However, it is possible for one of the
contour lines to pass exactly through the minimax point.  That contour
is a pair of crossed lines, as shown in the version on the right.  You
should compare these contour plots with the density plots of the same
functions on page~\pageref{density}, and with their graphs on
page~\pageref{saddle}.
	
\special{psfile=../ps/913/CONTOUR2.PS}  
\special{psfile=../ps/913/CONTOUR3.PS}  
\special{psfile=../ps/913/CONTOUR4.PS}  

 \vspace{120pt}

{\small
\begin{picture}(360,50)
	\put(52,12){\makebox(0,0){two functions but one plot}}
	\put(52,-2){\makebox(0,0){$z = x^2 + y^2$}}
	\put(52,-16){\makebox(0,0){$z = -x^2 - y^2$}}
	\put(283,8){\makebox(0,0){two plots of a single function}}
	\put(283,-8){\makebox(0,0){$z = x^2 - y^2$}}
\end{picture}
}

\newpage

There is a direct connection between the contour plot 
%
\mar{Contours are\\ horizontal slices\\ of a graph}
%
of a function and its graph.\index{graph!slice of
  a}\index{graph!horizontal slice} Contours are horizontal slices
of the graph, just as grid lines are vertical slices.  Below, we use
the standard functions $z = x^2 - y^2$ and $z = x^2 + y^2$ to
illustrate the connection.  Notice that every contour down in the $x,
y$-plane lies exactly below, and has the same shape as, a horizontal
slice of the graph.  This picture explains why contours are called
{\em level\/} curves.\index{level curve}
	
\special{psfile=../ps/913/CONSURF1.PS}
\special{psfile=../ps/913/CONSURF2.PS}
\vspace{190pt}
\label{standard min max-end}

To get some more experience with contour plots, we return to 
%
\mar{Energy of the pendulum, again}	
%
the energy\index{pendulum}\index{energy} function of the pendulum:
\[
E(\theta, v) = 1 - \cos{\theta} + \tfrac{1}{2} v^2.
\]
From the contour plot alone you should be able to see that $E$ has
either a minimum or a maximum at $(\theta, v) = (0, 0)$, and another
at $(2\pi, 0)$.  The contours also provide evidence that there is a
saddle (minimax) near $(\theta, v) = (-\pi, 0)$ and $(\pi, 0)$.  It is
also apparent that $E$ is a {\em periodic\/} function of $\theta$.

\special{psfile=../ps/913/contour5.ps}

\newpage

What you should find most striking about this plot, however, is the
way it resembles the phase portrait of the pendulum (chapter 8, pages
\pageref{pend-begin}--\pageref{pend-end}).
% qqq
Every level curve here looks like a trajectory\index{trajectory} of
the dynamical system.\index{dynamical system}
%
\mar{Energy contours\\ are trajectories\\ of the dynamics}	
%
This is no accident.  We know from chapter 8 that the energy is a
first integral for the dynamics.\index{first integral} In other words,
energy is constant along each trajectory---this is the law of
conservation of energy.  But each level curve shows where the energy
function has some fixed value.  Therefore, each trajectory must lie on
a single energy level.

\special{psfile=../ps/913/contour6.ps}

\noindent
\begin{minipage}[t]{3.3in}
\quad \ We can carry the connection between contours and trajectories
even further.  Closed trajectories correspond to {\em
  oscillations\/}\index{oscillation} of the pendulum.  But the closed
trajectories are the closed contours, and these are the ones that
surround the minimum.  In particular, they are {\em low\/} energy
levels.  By contrast, at higher energies ($E > 2$, in fact), the
pendulum will just continue to spin in what ever direction it was
moving initially.  Thus, each high energy level is occupied by {\em
  two\/} trajectories---one for clockwise spinning and one for
counter-clockwise.
\end{minipage}


\subsection{Technical Summary}

The examples we have seen so far were meant to introduce some of the
common ways of visualizing a function $z = f(x, y)$.  To use them most
effectively, though, you need to know more precisely how each is
defined.  We review here the definition of a graph, a density plot, a
contour plot, and a terraced density plot.
	
\special{psfile=../ps/914/SUMGRAPH.PS}

\medskip
\noindent
\begin{minipage}[t]{3.3in}
{\bf Graph}.\index{graph!surface} The graph of $z = f(x, y)$ lies in
the 3-dimensional space with coordinates $(x, y, z)$.  To construct
it, take any input $(x, y)$.  Identify this with the point $(x, y, 0)$
in the $x, y$-plane (which is defined by the condition $z = 0$).  The
corresponding point on the graph lies at the height $z = f(x, y)$
above the $x, y$-plane.  This point has coordinates $(x, y, f(x, y))$.
{\bf The graph is the set of all points of the form $(x, y, f(x,
  y))$.}  This is a 2-dimensional surface.
\end{minipage} 

\newpage

\special{psfile=../ps/914/SUMDENSE.PS}
\special{psfile=../ps/914/SUMCON2.PS}
\special{psfile=../ps/914/sumcon4.ps}
\hfill
\begin{minipage}[t]{3.7in}
{\bf Density plot}.\index{plot!density} The density plot of $z = f(x,
y)$ lies in the 2-dimensional $x, y$-plane.  Choose any rectangle
where the function is defined, and let $m$ and $M$ be the minimum and
maximum values, respectively, of $f(x, y)$ on the rectangle.  Define
\[
\rho(x, y) = \dfrac{f(x, y) - m}{M - m}.
\]
Then $\rho$ satisfies $0 \leq \rho(x, y) \leq 1$ on the rectangle; it
is called a {\bf density function} ($\rho$ is the Greek letter {\em
  rho\/}).  {\bf In the density plot, the density of ink---or
  darkness---at $(x, y)$ is $\rho(x, y)$.}
\smallskip

{\bf Contour plot}.\index{plot!contour}\index{contour} A {\bf contour}
of $z = f(x, y)$ is the set of points in the $x, y$-plane where $f$
has some fixed value:
\[
f(x, y) = c.
\]
That fixed value $c$ is called the {\bf level}\index{level} of the
contour.  (The two solid ovals in the figure at the left make up a
single contour.)  A contour is also called a level curve.  {\bf A
  contour plot of $f$ is a collection of curves $f(x, y) = c_j$ in the
  $x, y$-plane.}  In the plot it is customary to use constants $c_1,
c_2, \ldots$ that are equally spaced; that is, the interval between
one $c_j$ and the next always has the same value $\Delta c$.
\smallskip

\enlargethispage*{40pt} {\bf Terraced density
  plot}.\index{plot!density!terraced}\index{terraced density plot}
This is a contour plot in which the region between two adjacent
contours is shaded with ink of a single density.  If the contours are
at levels $c_1$ and $c_2$, then the density that is typically chosen
is the one for the level half-way between these two---that is, for
their average $(c_1 + c_2)/2$.  Each region is called a {\bf
  terrace}.\index{terrace} Often, a terraced density plot is drawn in
color, using different colors for each terrace.  Television weather
programs use terraced density plots to describe the temperature
forecast for a large region.
\end{minipage}

\aside{We find density plots everywhere.  A photograph is a density
  plot of the light that fell on the film when it was exposed.  A
  newspaper `half-tone' illustration is also a density plot of an
  image.}

\subsubsection{The pros and cons}

Each of these modes of visualization has advantages and disadvantages.
All are reasonably good at indicating the extremes (the maxima and
minima) of a function.  A contour plot needs some additional
information---for example, a label on each contour to indicate its
level---to distinguish between maxima and minima.  However, if you
want to know the numerical value of $f(x, y)$ at a particular point
$(x, y)$, a contour plot with labels offers more precision than a
density plot.  It's usually better than a graph, too.

Overall, a graph has the biggest visual impact, 
%
\mar{Plots are visually economical in comparison to graphs}	
%
but there is a cost.  It takes three dimensions to represent the graph
of a function of two variables, but only two to represent a plot.  The
cost is that extra dimension.  It means that we cannot draw the graph
of a function of three variables.  That would take four mutually
perpendicular axes---an impossibility in our three-dimensional space.
However, we can produce a contour plot.

\subsection{Contours of a Function of Three Variables}

We pause here for a brief glimpse of a large subject.  By analogy with
the definition for a function of two variables,
%
\mar{Contours and levels}	
%
 we say that a {\bf contour}\index{contour} of the function $f(x, y,
 z)$\index{function!of three variables} is the set of points $(x, y,
 z)$ that satisfy the equation
\[
f(x, y, z) = c,
\]
for some fixed number $c$.  We call $c$ the {\bf level}\index{level}
of the contour.

Let's find the contours of $w = x^2 + y^2 + z^2$.  
%
\mar{The standard minimum}	
%
This is completely analogous to the function $x^2 + y^2$ with two
inputs.  (What do the contours of $x^2 + y^2$ look like?)  In
particular, $w$ has a minimum when $(x, y, z) = (0, 0,
0)$.\index{minimum} As the following diagram shows, $w = x^2 + y^2 +
z^2$ is the square of the distance from the origin to $(x, y, z)$.
(We use the Pythagorean theorem twice: once for $p^2$ and once for
$q^2$.)  Consequently, all points $(x, y, z)$ where $w$ has a fixed
value \linebreak

\special{psfile=../ps/915/XYZ1.PS}

\begin{align*}
\rule{200pt}{0pt}	p^2 &= x^2 + y^2, \\[3pt]
	q^2 &= p^2 + z^2 \\
	    &= x^2 + y^2 + z^2 \\
	    &= w.
	\end{align*}
	
\newpage

\noindent
lie a fixed distance from the origin.  Specifically, $w = x^2 + y^2 +
z^2 = c$ is the following set:
\begin{itemize}

\item $c > 0$: the sphere of radius $\sqrt{c}$ entered at the origin;

\item $c = 0$: the origin itself;

\item $c < 0$: the empty set.

\end{itemize}
The contour plot of $w = x^2 + y^2 + z^2$ is thus 
%
\mar{The contours \\are spheres}
%
a nest of concentric spheres, as shown in the illustration
below.\index{level surface} The value of $w$ is constant on each
sphere.  (The tops of the spheres have been cut away so you can see
how the spheres nest; the whole thing resembles an onion.)

\special{psfile=../ps/915/SHELLS1.PS}

\newpage

Below is the contour plot of another standard function 
%
\mar{A standard minimax}	
%
with three input variables:
\[
w = f(x, y, z) = x^2 + y^2 - z^2.
\]
A quarter of each surface has been cut away so you can see how the
surfaces nest together.  Note that $w = 0$ is a cone, and every
surface with $w < 0$ consists of two disconnected (but congruent)
pieces---an upper half and a lower half.\index{minimax}

You should compare this function to the standard minimax $x^2 - y^2$
in two variables.  The three-variable function $w$ has a minimum with
respect to {\em both\/} of the variables $x$ and $y$, while it has a
maximum with respect to $z$.  (Do you see why?  The arguments are
exactly the same as they were for two input variables on
page~\pageref{saddle}.)
%
\mar{The contours \\are hyperbolic shapes}	
%
Furthermore, the contours of $x^2 - y^2$ are a family of hyperbolas,
and the contours of $x^2 + y^2 - z^2$ are surfaces obtained by
rotating these hyperbolas about a common axis.

\special{psfile=../ps/915/SHELLS2.PS}
\special{psfile=../ps/915/SHELLS3.PS}

\newpage

It is a general fact---and our two examples 
%
\mar{When there are three input variables, the contours are surfaces}	
%
provide good evidence for it---that a single contour of a function of
three variables is a {\em surface}.  Thus a contour is a curve or a
surface, depending on the number of input variables.  We often use the
term {\bf level set}\index{level set} (rather than a level {\em
  curve\/} or a level {\em surface\/})\index{level
  surface}\index{surface!level} as a generic name for a contour.

\subsection{Exercises}

In many of these exercises it will be essential to have a computer
program to make graphs, terraced density plots, and contour plots of
functions of two variables.
	
\numa Use a computer to obtain a graph of the function $z = \sin{x}
\sin{y}$ on the domain $0 \leq x \leq 2 \pi$, $0 \leq y \leq 2 \pi$.
How many maximum points do you see?  How many minimum points?  How
many saddles?

\sub Determine, as well as you can from the graph, the location of the
maximum, minimum, and saddle points.

\num Continuation.  Make the domain $-2 \pi \leq x \leq 2 \pi$, $-2
\pi \leq y \leq 2 \pi$ and answer the same questions you did in the
previous exercise.  (Does the graph look like an egg carton?)

\num Obtain a terraced density plot (or a contour plot) of $z =
\sin{x} \sin{y}$ on the domain $-2 \pi \leq x \leq 2 \pi$, $-2 \pi
\leq y \leq 2 \pi$.  Locate the maximum, minimum, and saddle points of
the function.  Do these results agree with those from the previous
exercise?

\num Obtain the graph of $z = \sin{x} \cos{y}$ on the domain $0 \leq x
\leq 2 \pi$, $0 \leq y \leq 2 \pi$.  How does this graph differ from
the one in exercise~1?  In what ways is it similar?

\num\label{quartic} Obtain the graph of $z = 2x + 4 x^2 - x^4 - y^2$
when $-2 \leq x \leq 2$, $-4 \leq y \leq 4$.  Locate all the minimum,
maximum, and saddle points in this domain.  [Note: the minimum is on
  the boundary!]

\num Continuation.  Obtain a terraced density plot (or contour plot)
for the function in the previous exercise, using the same domain.  Use
the plot to locate all the minimum, maximum, and saddle points.
Compare your results with those of the previous exercise.

\numa Obtain the graph of $z = 2x - y$ on the domain $-2 \leq x \leq
2$, $-2 \leq y \leq 2$.  What is the shape of the graph?

\sub Graph the same function of the domain $2 \leq x \leq 6$, $0 \leq
y \leq 4$.  What is the shape of the graph?  How does this graph
compare to the one in part (a)?

\numa Continuation.  Sketch three different slices of the graph of $z
= 2x - y$ in the $y$-direction.  What do the slices have in common?
How are they different?

\sub Answer the same questions for slices in the $x$-direction.

\numa Obtain the graph of $z = .3 x + .8 y + 2.3$; choose the domain
yourself.  Where does the graph intercept the $z$-axis?

\sub Describe the vertical slices of this graph in the $y$-direction
and in the $x$-direction.

\num Describe the vertical slices of the graph of $z = px + qy + r$ in
the $y$-direction and in the $x$-direction.

\numa Compare the contours of the function $z = x^2 + 2 y^2$ to those
of $z = x^2 + y^2$.

\sub What is the shape of the graph of $z = x^2 + 2 y^2$?  Decide this
first using only the information you have about the contours.  Then
use a computer to obtain the graph.

\numa Compare the contours of the function $z = x^2 - 2 y^2$ to those
of $z = x^2 - y^2$.

\sub What is the shape of the graph of $z = x^2 - 2 y^2$?  Decide this
first using only the information you have about the contours.  Then
use a computer to obtain the graph.

\numa Obtain a contour plot of the function $z = x^2 + xy + y^2$.

\sub What is the shape of the graph of $z = x^2 + xy + y^2$?  Decide
this first using only the information you have about the contours.
Then use a computer to obtain the graph.

\numa Obtain a contour plot of the function $z = x^2 + 3xy + y^2$.

\sub What is the shape of the graph of $z = x^2 + 3xy + y^2$?  Decide
this first using only the information you have about the contours.
Then use a computer to obtain the graph.

\numa Obtain a contour plot of the function $z = x^2 + 2xy + y^2$.

\sub What is the shape of the graph of $z = x^2 + 2xy + y^2$?  Decide
this first using only the information you have about the contours.
Then use a computer to obtain the graph.

\num Complete this statement: The function $f(x, y; p) = x^2 + p xy +
4 y^2$, which depends on the parameter $p$, has a minimum at the
origin when \rule[-1pt]{30pt}{.3pt} and a minimax when
\rule[-1pt]{30pt}{.3pt}\,.

\numa Obtain the graph and a terraced density plot of the function $z
= 3x^2 + 17 xy + 12 y^2$.  What is the shape of the graph?

\sub What is the shape of the contours?  Indicate how the contours fit
on the graph.

\numa Obtain the graph and a terraced density plot of the function $z
= 3x^2 + 7 xy + 12 y^2$.  What is the shape of the graph?

\sub What is the shape of the contours?  Indicate how the contours fit
on the graph.

\numa Obtain the graph and a terraced density plot of the function $z
= 3x^2 + 12 xy + 12 y^2$.  What is the shape of the graph?

\sub What is the shape of the contours?  Indicate how the contours fit
on the graph.


\num Obtain the graph of $z = f(x, y) = xy$ on the domain $-3 \leq x
\leq 3$, $-3 \leq y \leq 3$.  Does this function have a maximum or a
minimum or a saddle point?  Where?

\numa Continuation.  Sketch slices of the graph of $z = xy$ in the
$y$-direction, for each of the values $x = -2$, $-1$, 0, 1, and 2.
What is the general shape of each of these slices?

\sub Repeat part (a), but make the five slices in the
$x$-direction---that is, fix $y$ instead of $x$.

\num Continuation.  Show how the slices you obtained in the previous
exercise fit (or appear) on the graph you obtained in the exercise
just before that one.

\numa Continuation.  Let $x = u + v$ and $y = u - v$.  Express $z$ in
terms of $u$ and $v$, using the fact that $z = xy$.  Then obtain the
graph of $z$ as a function of the new variables $u$ and $v$.

\sub What is the shape of the graph you just obtained?  Compare it to
the graph of $z = xy$ you obtained earlier.

\num Can you draw a network of straight lines on the saddle surface $z
= x^2 - y^2$?

\num Obtain a terraced density plot of $z = xy$.  How do the contours
of this plot fit on the graph of $z = xy$ you obtained in a previous
exercise?

\num The graphs of $z = x^2 + 5 xy + 10 y^2$ and $z = 3$ intersect in
a curve.  What is the shape of that curve?

\num The graphs of $z = x^2 + y^3$ and $z = 0$ intersect in a curve.
What is the shape of that curve?

\num The graphs of $z = 2x - y$ and $z = .3 x + .8 y + 2.3$ intersect
in a curve.  What is the shape of that curve?


\subsubsection{First integrals}
\vspace{-10pt}
\index{first integral}

\num A hard spring\index{hard spring}\index{spring!hard} described by
the dynamical system
\[
\dfrac{dx}{dt} = v, \qquad \quad \dfrac{dv}{dt} = - cx - \beta x^3,
\]
has a first integral of the form
\[
E(x, v) = \tfrac {1}{2}c x^2 + \tfrac {1}{2} \beta x^4 + \tfrac{1}{2}v^2.
\]
This is the {\bf energy}\index{energy} of the spring.  (See
chapter~7.3, especially exercise~13, page~\pageref{spring energy}.)

\sub Let $c = 16$ and $\beta = 1$.  Obtain the graph of $E(x, v)$ on a
domain that has the origin at its center.  Locate all the minimum,
maximum, and saddle points in this domain.

\sub What is the state of the spring (that is, its position $x$ and
its velocity $v$) when it has minimum energy?

\num A soft spring\index{soft spring}\index{spring!soft} described by
the dynamical system
\[
\dfrac{dx}{dt} = v, \qquad \quad \dfrac{dv}{dt} = \dfrac{-25x}{1 + x^2},
\]
has an energy integral of the form
\[
E(x, v) = \tfrac{25}{2} \ln(1 + x^2) + \tfrac{1}{2}v^2.
\]
(See exercise~16, page~\pageref{ln first int}.)  

\sub Obtain the graph of $E(x, v)$.  Experiment with different
possibilities for the domain until you get a good representation.

\sub Obtain a terraced density plot of $E(x, v)$ over the same domain
you chose in part (a).  Compare the two representations of $E$.

\sub Does the spring have a state of minimum energy?  If so, where is
it?

\sub Does the spring have a state of {\em maximum\/} energy?  Explain
your answer.

\numa {\bf The Lotka--Volterra equations}.\index{Lotka--Volterra
  equations} According to exercise~33 of chapter~7.3
(page~\pageref{L-V first int}), the function
\[
E(x, y) = .1 \ln{y} + .04 \ln{x} - .005 \, y -.004 \, x
\]
is a first integral of the dynamical system
\begin{align*}
x' &= .1 x - .005 \, xy,	\\
y' &= .004 \, xy - .04 \, y.
\end{align*}
Obtain the graph of $E$ on the domain $1\leq x \leq 50$, $1 \leq y
\leq 50$.  (Why not enlarge the domain to $0 \leq x \leq 50$, $0 \leq
y \leq 50$?)

\sub Find all maximum, minimum, and saddle points on this graph.  What
is the connection between the maximum of $E$ and the equilibrium point
of the dynamical system?

\sub Obtain a contour plot of $E$ on the same domain as in part (a).
Compare the contours of $E$ and the trajectories of the dynamical
system.  (This reveals a {\em conservation of ``energy''\/} for the
solutions of the Lotka-Volterra equations.)

\numa Continuation.  Here is another first integral of the same
dynamical system as in the previous exercise:
\[
H(x, y) = \dfrac{x^{0.04} y^{0.1}}{\exp(.005 \, y + .004 \, x)}.
\]
Obtain the graph of $H$ and compare it to the graph of $E$ in the
previous exercise.

\sub Obtain a contour plot of $H$, and compare the contours to the
trajectories of the dynamical system.

\newpage
%%%%%%%%%%%%%%%%%%%%%%%%%%%%%%%%%%%%%%%%%%%%%%%%%%%%%%%%%%%%%%%%%%%%%%%%%%
%                                                                        %
%                                Section 2                               %
%                                                                        %
%%%%%%%%%%%%%%%%%%%%%%%%%%%%%%%%%%%%%%%%%%%%%%%%%%%%%%%%%%%%%%%%%%%%%%%%%%

\section{Local Linearity}
\index{local linearity}

Local linearity is the central idea of chapter 3: it says that a graph
looks straight when viewed under a microscope.\index{microscope} Using
this observation we were able to give a precise meaning to the {\em
  rate of change\/} of a function and, as a consequence, to see why
Euler's method produces solutions to differential equations.  At the
time we concentrated on functions with a single input variable.  In
this section we explore local linearity for functions with two or more
input variables.

\subsection{Microscopic Views}

Consider the cubic $f(x, y) = x^3 - 4x - y^2$ that we used as an
example in the previous section.  We'll examine both the graph and the
plot of $f$ under a microscope.\index{graph!magnifying
  a}\index{graph}\index{graph!surface}
%
\mar{Magnifying \\a graph}	
%
In the figure below we see successive magnifications of the graph near
the point where $(x, y) = (1.5, -1)$.  The initial graph, in the left
rear, is drawn over the square
\[
-2.5 \leq x \leq 2.5, \qquad -2.5 \leq y \leq 2.5.
\]
With each magnification, the portion of the surface we see bends less
and less.  {\bf The graph approaches the shape of a flat plane}.

\special{psfile=../ps/921/micro1.ps}
\special{psfile=../ps/921/micro2.ps}
\special{psfile=../ps/921/micro3.ps}

\newpage
	
Contour plots\index{plot!contour} for $f(x, y) = x^3 - 4x - y^2$
appear below.  Again, we magnify near the point $(x, y) = (1.5, -1)$.
Each window below is a small part of the window to its left.  In the
large scale plot, which is the first one on the left, the contours are
quite variable in their direction and spacing.  With each
magnification, that variability decreases.  {\bf The contours become
  straight, parallel, and equally
  spaced}.\index{contour}\label{plots1}

\special{psfile=../ps/921/micro4.ps}
\vspace{158pt}

The process of magnification thus leads us to functions whose graphs
are flat and whose contours are straight, parallel, and
equally-spaced.  As we shall now see, these are the linear functions.

\subsection{Linear Functions}
\index{linear function}
\index{function!linear}

A linear function 
%
\mar{Responses to\\ changes in input}	
%
is defined by the way its output responds to {\em changes\/} in the
input.  Specifically, in chapter 1 we said
\[
y = f(x) \quad \mbox{is linear if} \quad \Delta y = m \cdot \Delta x.
\]
This is the simplest possibility: changes in output are strictly
proportional to changes in input.  The multiplier\index{multiplier}
$m$ is both the {\em rate\/} at which $y$ changes with respect to $x$
and the {\em slope\/} of the graph of $f$.

Exactly same idea defines a linear function 
%
\mar{The definition}	
%
of two or more variables: the change in output is strictly
proportional to the change in any one of the inputs.
\begin{center}
\begin{picture}(330,70)		
\put(0,0) {
  \framebox(330,70){
    \begin{minipage}{300pt}			
      {\bf Definition}. The function $z = f(x_1, x_2, \ldots, x_n)$ is 
      {\bf linear} \\if there are multipliers $p_1, p_2, \ldots, p_n$ for which
      \vspace{-20pt}
      
	$$
	\Delta z = p_1 \cdot \Delta x_1, \quad \Delta z = p_2 \cdot \Delta x_2,
		\quad \ldots, \quad \Delta z = p_n \cdot \Delta x_n.
	$$
    \end{minipage}
  }
}
\end{picture}
\end{center}
There is one multiplier for each input variable.  The multipliers are
constants and they are, in general, all different.

The definition describes how $z$ responds 
%
\mar{Partial and \\total changes}	
%
to each input separately.  We call each $p_j \cdot \Delta x_j$ a {\bf
  partial change}.\index{change!partial} The multiplier $p_j = \Delta
z/\Delta x_j$ is the corresponding {\bf partial rate of
  change}.\index{partial rate of change}\index{rate of change!partial}
Of course, several input variables may change simultaneously.  In that
case, the {\bf total change}\index{total change}\index{change!total}
in $z$ will just be the sum of the individual changes produced by the
several variables:
\[
\Delta z = p_1 \cdot \Delta x_1 + p_2 \cdot \Delta x_2 +
	 	\cdots + p_n \cdot \Delta x_n.
\]

Of course, if the total change of a function satisfies 
%
\mar{Another way \\to describe a \\linear function}	
%
this condition, then each partial change has the form $p_j \cdot
\Delta x_j$.  (If only $x_j$ changes, then all the other $\Delta x_k$
must be 0.  So $\Delta z$ becomes simply $p_j \cdot \Delta x_j$.)
Consequently, the function must be linear.  In other words, we can use
the formula for the total change as another way to define a linear
function.
\label{altdef}% 
\begin{center}
\begin{picture}(330,70)		
\put(0,0) {
  \framebox(330,70){
    \begin{minipage}{300pt}			
      {\bf Alternate definition}. The function $z = f(x_1, x_2, \ldots, x_n)$\\
      is {\bf linear} if there are multipliers $p_1, p_2, \ldots, p_n$ for 
      which
      \vspace{-10pt}
        $$
	\Delta z = p_1 \cdot \Delta x_1 + p_2 \cdot \Delta x_2 +
	 	\cdots + p_n \cdot \Delta x_n.
	$$
    \end{minipage}
  }
}
\end{picture}
\end{center}

\subsubsection{Formulas for linear functions}  

When $z = f(x_1, x_2, \ldots, x_n)$ is a linear function, 
%
\mar{From the definition\\ to a formula}	
%
we know how $\Delta z$ depends on the changes $\Delta x_j$, but that
doesn't tell us explicitly how $z$ itself is related to the input
variables $x_j$.  There are several ways to express this relation as a
formula,\index{linear function!formula} depending on the nature of
the information we have about the function.  For the sake of clarity,
we'll develop these formulas first for a function of two variables: $z
= f(x, y)$.
\smallskip

\noindent
$\bullet$ {\bf The initial-value form}.\index{linear function!initial-value form}\index{initial-value form}
%
\mar{Given the partial\\ rates of change and\\ an initial point}
%
Suppose we know the value of a linear function at some given
point---called the {\bf initial point}---and\index{initial point} we
also know its partial rates of change.  Can we construct a formula for
the function?  Suppose $z = z_0$ when $(x, y) = (x_0, y_0)$, and
suppose the partial rates of change are
\[
p = \dfrac{\Delta z}{\Delta x} \qquad \mbox{and} \qquad
		q = \dfrac{\Delta z}{\Delta y}.
\]
If we let
\[
\Delta x = x - x_0, \qquad \Delta y = y - y_0, \qquad \Delta z = z - z_0,
\]
then we can write
\begin{align*}
z - z_0 = \Delta z &= p \cdot \Delta x + q \cdot \Delta y \\
		   &= p \cdot (x - x_0) + q \cdot (y - y_0).
\end{align*}
This is the initial-value form of a linear function.  For example, if
the initial point is $(x, y) = (4, 3)$, $z = 5$, and the partial rates
of change are $\Delta z/\Delta x = -\tfrac{1}{2}$, $\Delta z/\Delta y
= +1$, the equation of the linear function can be written
\[
z - 5 = -\tfrac{1}{2} (x - 4) + (y - 3).
\]

\smallskip
\noindent
$\bullet$ {\bf The intercept form}.\index{linear function!intercept
  form}\index{intercept form}
%
\mar{Given the partial\\ rates of change and\\ the $z$-intercept}
%
This is a special case of the initial-value form, in which the initial
point is the origin: $(x, y) = (0, 0)$, $z = r$.  The formula becomes
\[
z - r = p  x + q y, \qquad \mbox{or} \qquad
		z = p x + q  y + r.
\]
As we shall see, the graph of this function in $x, y, z$-space passes
through the point $(x, y, z) = (0, 0, r)$ on the $z$-axis.  This point
is called the \mbox{\bf{\boldmath
    $z$}-intercept}\index{intercept}\index{$z$-intercept} of the
graph.  Sometimes we simply call the number $r$ itself the
$z$-intercept.

Notice that, with a little algebra, we can convert the previous
example to the form $z = -\tfrac{1}{2}x + y + 4$.  This is the
intercept form, and the $z$-intercept is $z = 4$.
\smallskip

If there are $n$ input variables, $x_1$, $x_2$, \ldots , $x_n$, 
%
\mar{\vspace{44pt}The form of\\ a linear function\\ of $n$ variables}	
%
instead of two, and an initial point has coordiantes $x_1^0$, $x_2^0$,
\ldots, $x_n^0$, then a linear equation has the following forms:
\begin{alignat*}{2}
\textbf{initial-value:}&\quad &  z - z_0 
      &= p_1 (x_1 - x_1^0) + p_2 (x_2 - x_2^0) 
		+ \cdots + p_n (x_n - x_n^0), \\
\textbf{\boldmath $z$-intercept:}& &  z &= p_1  x_1 + p_2  x_2 
		+ \cdots + p_n  x_n + r.
\end{alignat*}

\subsubsection{The graph of a linear function}
\index{graph!linear function}
\index{linear function!graph}

On the left at the top of the next page is the graph of the linear
function
\[
z = \tfrac{1}{2} x  + y + 4.
\]
The graph is a flat plane.  In particular, grid lines\index{grid line}
parallel to the $x$-axis (which represent vertical slices\index{slice}
with $y = c$) are all straight lines with the same slope $\Delta
z/\Delta x = -\tfrac{1}{2}$.  The other grid lines (with $x = c$) are
all straight lines with the same slope $\Delta z/\Delta y = +1$.

\newpage

\mbox{}
 \special{psfile=../ps/922/linear1.ps}
 \special{psfile=../ps/922/linear2.ps}
\vspace{144pt}

{\small
	\noindent
	\hspace{40pt}
	$z = - \tfrac{1}{2} x  + y + 4$
	\hspace{135pt}
	$z = px + qy + r$
}
\vspace{10pt}

On the right, above, is the graph of the general linear function
written in intercept form: $z = px + qy + r$.  The graph is the plane
that can be identified by three distinguishing features:

$\bullet$ it has slope\index{slope} $p$ in the $x$-direction; 
	
$\bullet$ it has slope $q$ in the $y$-direction.

$\bullet$ it intercepts the $z$-axis at $z = r$; 
\smallskip
	
Let's see how we can {\em deduce\/} 	 
%
\mar{The definition of\\ a linear function\\ implies that its graph\\ is a
  flat plane}
%
that the graph must be this plane.  First of all, the partial rate
$\Delta z/\Delta x$\index{rate of change!partial} tells us how $z$
changes when $y$ is held fixed.  But if we fix $y = c$, we get a
vertical slice of the graph in the $x$-direction.  The slope if that
vertical slice is $\Delta z/\Delta x = p$.  Since $p$ is constant, the
slice is a straight line.  The value of $y = c$ determines which slice
we are looking at.  Since $\Delta z/\Delta x$ doesn't depend on $y$,
all the slices in the $x$-direction have the {\em same\/} slope.
Similarly, all the slices in the $y$-direction are straight lines with
the same slope $q$.  The only surface that can be covered by a grid of
straight lines in this way is a flat plane.  Finally, since $z = r$
when $(x, y) = (0, 0)$, the graph intercepts the $z$-axis at $z = r$.

\subsubsection{Contours of a linear function}
\index{linear function!contours}
\index{contour!linear function}

A contour of {\em any\/} function 
%
\mar{Each contour is\\ a straight line}	
%
$f(x, y)$ is the set of points in the $x, y$-plane where $f(x, y) =
c$, for some given constant $c$.  If $f = px + qy + r$, then a contour
has the equation
\[
px + qy + r = c \qquad \mbox{or} \qquad 
		y = - \dfrac{p}{q}x + \dfrac{c - r}{q}.
\]
This is an ordinary straight line in the $x, y$-plane.  Its slope is
$-p/q$ and its $y$-intercept is $(c - r)/q$.  \label{contourslope} (If
$q = 0$ we can't do these divisions.  However, this causes no problem;
you should check that the contour is just the vertical line $x = (c -
r)/p$.)

\special{psfile=../ps/922/linear3.ps}

\noindent
\begin{minipage}[t]{3.5in}
\quad \ To construct a contour plot, we must give the constant $c$ a
sequence of equally-spaced values $c_j$, with $c_{j+1} = c_j + \Delta
c$.  This generates a sequence of straight lines
\[
px + qy + r = c_j, \quad \mbox{or} \quad 
		y = - \dfrac{p}{q}x + \dfrac{c_j - r}{q}.
\]
These lines all have the same slope $-p/q$, so they are parallel.
(Notice the value of $c$ doesn't affect the slope.)  The $y$-intercept
of the $j$-th contour is $(c_j - r)/q$.  Therefore, the distance along
the $y$-axis \linebreak
\end{minipage}

\vspace{-8pt}

\noindent
between one intercept and the next is $\Delta c/q$.  The contours are
thus straight, parallel, and equally-spaced.  (You should check that
this is still true if $q = 0$.)  Note that the figure at the left,
above, is drawn with $\Delta c > 0$ but $q < 0$.

\subsubsection{Geometric interpretation of the partial rates}

What happens to the graph or the contour plot if you double one of the
partial rates of change\index{rate of change!partial}\index{partial
  slope}\index{slope!partial} of a linear function?
%
\mar{Partial rates and partial slopes}	
%
The graph on the right, below, shows the effect of doubling the
partial rate with respect to $x$ of the function $z = - \tfrac{1}{2}x
+ y + 4$.  As you can see, the slope in the $x$-direction \linebreak

\vspace{-5pt}
	
\special{psfile=../ps/922/linear4.ps}
\special{psfile=../ps/922/linear5.ps}
\vspace{160pt}

{\small
	\noindent
	\hspace{55pt}
	$z = -\tfrac{1}{2} x  + y + 4$
	\hspace{105pt}
	$z = -x + y + 4$
}

\newpage

\noindent
is doubled (from $-\tfrac{1}{2}$ to $-1$).  Had we increased the
partial rate by a factor of~10, the slope would have increased by a
factor of~10 as well.  Notice that the slope in the $y$-direction is
not affected.  
%
\mar{The overall tilt\\ of a graph is altered}
%
Nevertheless, the overall `tilt' of the graph {\em has\/} been
altered.  We shall have more to say about this feature in a moment,
when we introduce the {\bf gradient}\index{gradient} of a linear
function to describe the overall tilt.
\medskip

A change in the partial rates has a more complex effect on the contour
plot.  Perhaps it is more surprising, too.  To make valid comparisons,
we have constructed all three plots below with the same spacing
between levels (namely $\Delta z = 1$).  Notice how the levels meet
the $x$- and $y$-axes in the plot on the left ($z = -tfrac{1}{2}x + y
+4$).
%
\mar{Partial rate and the spacing of contours} 	
%
For each unit step we take along the $y$-axis, the $z$-value increases
by 1.  This is the meaning of $\Delta z/\Delta y = +1$.  By contrast,
we have to take {\em two\/} unit steps along the $x$-axis to produce
the same size change in $z$.  Moreover, $z$ {\em decreases\/} by 1
when $x$ increases by 2.  This is the meaning of $\Delta z/\Delta x =
-\tfrac{1}{2}$.  In particular, the relatively wide spacing between
$z$-levels along the $x$-axis reflects the relative smallness of
$\Delta z/\Delta x$.
 


\vspace{13pt}
\special{psfile=../ps/922/linear6.ps}
\special{psfile=../ps/922/linear7.ps}
\special{psfile=../ps/922/linear8.ps}
\vspace{138pt}

{\small
	\noindent
	\hspace{-3pt}
	$z = -\tfrac{1}{2} x  + y + 4
	\hspace{65pt}
	z = -x + y + 4$
	\hspace{65pt}
	\mbox{$z = -x + 2y + 4$}
}

\vspace{5pt}

Therefore, when we double the size of $\Delta z/\Delta x$---as we do
in the middle plot---we should cut in half
%
\mar{The larger \\the partial rate, \\the closer \\the contours} 
%
the spacing between $z$-levels along the $x$-axis.  As you can see,
this is exactly what happens.  Notice that the spacing along the
$y$-axis is not altered.  Consequently, the contours change direction
and they get packed more closely together.

Suppose we double {\em both\/} partial rates---as we do in the plot on
the right.  Then the spacing between contours is cut in half along
both axes.  Because the change is uniform, the contours keep their
original direction.

\newpage

\subsection{The Gradient of a Linear Function}
\index{linear function!gradient}
\index{gradient}

By making use of the concept of a vector,\index{vector} introduced in
the last chapter, we can construct
%
\mar{The vector of\\ partial rates}	
%
still another geometric interpretation of the partial
rates\index{rate of change!partial} of a linear function. This vector
is called the gradient, and it is defined in the following way.
\begin{center}
\begin{picture}(340,75)		
\put(0,0) {
  \framebox(340,75){
    \begin{minipage}{310pt}			
{\bf Definition}.  The {\bf gradient} of a linear function $z = f(x,
y)$ is the vector whose components are its partial rates of
change:
\linebreak%
\vspace{-8pt}
$$
\mbox{grad $z$} = \nabla z\index{$\nabla$}  
  = \left(\dfrac{\Delta z}{\Delta x}, \dfrac{\Delta z}{\Delta y}\right).
$$ 
    \end{minipage}
  }
}
\end{picture}
\end{center}
	
The gradient is perhaps the most concise and useful tool for
describing the growth of a
%
\mar{The direction of\\ most rapid growth}	
%
function of several variables.\index{growth!most rapid} To get an
idea of the role that it plays, consider this question: {\em In what
  direction should we move from a given point in the $x, y$-plane so
  that the value of a linear function increases most rapidly\/}?

\vspace{10pt}
	
\special{psfile=../ps/922/linear9.ps}
\special{psfile=../ps/922/linear10.ps}
\special{psfile=../ps/922/linear11.ps}

Of course, the answer will depend on the linear function.  
%
\mar{\hspace{-20pt} $z = -\tfrac{1}{2} x + y + 4$}		
%
Let's use $z = -\tfrac{1}{2} x + y + 4$ and start from the point $(x,
y) = (2.4, 1.6)$.  We can make $z$ undergo a very large change simply
by moving very far from this point.  Therefore, to make valid
comparisons, we will restrict ourselves to motions that carry us
exactly one unit of distance in various directions.  The vectors in
the figure at the right show some of the possibilities.  Their tips lie
on a circle of radius~1.
	
Thus, to choose the direction in which $z$ increases most rapidly, we
must simply find the point on this circle where the value of $z$ is
largest.  The contour line at this level must be tangent to the
circle.  The vector perpendicular to this contour line (see the second
figure) therefore points in the direction of most rapid growth.  Since
perpendiculars have negative reciprocal slopes, and since all the
contour lines have slope $+1/2$, it follows that the vector must have
slope $-2/1$.
	
At the left is a magnified view of this vector.  We know 
\[
\Delta x < 0, \qquad \dfrac{\Delta y}{\Delta x} = -2, \qquad 
	\mbox{and} \qquad  (\Delta x)^2 + (\Delta y)^2 = 1.
\]
Thus $\Delta y = -2 \cdot \Delta x$, so $(\Delta x)^2 + 4 \, (\Delta
x)^2 = 1$.  This implies
\[
5 \, (\Delta x)^2 = 1, \qquad \mbox{so} \qquad 
	\Delta x = \dfrac{-1}{\sqrt{5}}, \quad 
	\Delta y = \dfrac{2}{\sqrt{5}}.
\]
	
Thus, among all the motions $(\Delta x, \Delta y)$ we have considered, we obtain the greatest change in $z$ by choosing
\[
(\Delta x, \Delta y) = 
		\left(\dfrac{-1}{\sqrt{5}}, \dfrac{2}{\sqrt{5}}\right).
\]
To determine how large this change is, we can use 
%
\mar{The magnitude of\\ most rapid growth}
%
the alternate definition of a linear function (see
page~\pageref{altdef})
	$$
	\Delta z = \dfrac{\Delta z}{\Delta x} \cdot \Delta x + 
			\dfrac{\Delta z}{\Delta y} \cdot \Delta y
		= -\dfrac{1}{2} \cdot \dfrac{-1}{\sqrt{5}} + 
			1 \cdot \dfrac{2}{\sqrt{5}}
		= \dfrac{5}{2 \sqrt{5}} = \dfrac{\sqrt{5}}{2}.
	$$

\special{psfile=../ps/922/linear12.ps}
	
The gradient vector quickly gives us all this information.  First of all, the gradient vector has the value
\[
\mbox{grad $z$} = 
	\left( \dfrac{\Delta z}{\Delta x}, \dfrac{\Delta z}{\Delta y}\right) 
	= \left( -\tfrac{1}{2}, 1 \right).
\]
Since its slope is $1/-\tfrac{1}{2} = -2$, we see that it does indeed
point in the direction of most rapid growth.
%	
\mar{Information from\\ the gradient}	
%
Consequently, it is also perpendicular to the contour
line. Furthermore, its {\em length\/} gives the maximum growth rate.
We can see this by calculating the length using the Pythagorean
theorem:\index{Pythagorean theorem}
	$$
	\mbox{length} = \sqrt{\left(\dfrac{\Delta z}{\Delta x}\right)^2 +
	 	\left(\dfrac{\Delta z}{\Delta y} \right)^2} =
		\sqrt{\dfrac{1}{4} + 1} = 
		\sqrt{\dfrac{5}{4}} = 
		\dfrac{\sqrt{5}}{2}.
	$$
Our findings with this example point to the following conclusion.
	
\begin{center}
\begin{picture}(360,70)(5.2,0)		
\put(0,0) {
  \framebox(360,70){
    \begin{minipage}{330pt}			
      {\bf Theorem}.  The gradient of the linear function $z = px + qy + r$ is perpendicular to its contour lines.  It points in the direction in which $z$ increases most rapidly, and its length is equal to the maximum rate of increase.\label{gradtheorem} 
    \end{minipage}
  }
}
\end{picture}
\end{center}

Let's see why this is true.   
%
\mar{A proof}	
%
According to the observation on the previous page, the direction of
most rapid increase will be perpendicular to the contour lines.  The
gradient of $z = px + qy + r$ is the vector $\nabla z = (p, q)$.  Its
slope is $q/p$.  On page~\pageref{contourslope} we saw that the slope
of the contour lines is $-p/q$.  Since these slopes are negative
reciprocals, the gradient is indeed perpendicular to the contour
lines.
\label{grad}   

To determine the maximum rate of increase, we must see how much $z$
increases when we move exactly 1 unit of distance in the gradient
direction.  The gradient vector is $(p, q)$, and its length is
$\sqrt{p^2 + q^2}$.  Therefore, the vector
	$$
	(\Delta x, \Delta y) = 
	\left(\dfrac{p}{\sqrt{p^2 + q^2}}, \dfrac{q}{\sqrt{p^2 + q^2}}\right)
	$$
is 1 unit long and in the same direction as the gradient.  The
increase in $z$ along this vector is
	$$
\Delta z = p \cdot \Delta x + q \cdot \Delta y =
   p \cdot \dfrac{p}{\sqrt{p^2 + q^2}} + q \cdot \dfrac{q}{\sqrt{p^2 + q^2}}
   = \dfrac{p^2 + q^2}{\sqrt{p^2 + q^2}} = \sqrt{p^2 + q^2}.
	$$
This {\em is\/} the length of the gradient vector, 
%
\mar{End of the proof}
%
so we have confirmed that the length of the gradient is equal to the
maximum rate of increase.
	
\vspace{13pt}	
\special{psfile=../ps/922/linear13.ps}
\special{psfile=../ps/922/linear14.ps}
\special{psfile=../ps/922/linear15.ps}
\vspace{138pt}

{\small
	\noindent
	\hspace{20pt}
	$z = -\tfrac{1}{2} x  + y + 4$
	\hspace{65pt}
	$z = -x + y + 4$
	\hspace{65pt}
	\mbox{$z = -x + 2y + 4$}
}

\vspace{6pt} 

Shown above are the three linear functions we've already
examined.\index{plot!contour} In each case the gradient vector is
perpendicular to the contours,
%
\mar{Contour spacing and\\ the length \\of the gradient}
%
and it gets longer as the space between the contours {\em decreases}.
This is to be expected because the space between contours is also an
indicator of the maximum rate of growth of the function.
Widely-spaced contours tell us that $z$ changes relatively little as
$x$ and $y$ change; closely-spaced contours tell us that $z$ changes a
lot as $x$ and $y$ change.
\medskip
	
The connection between the\index{gradient!geometric interpretation} 
%
\mar{The gradient points directly uphill}	
%
gradient and the graph is particularly simple. Since the gradient
(which is a vector in the $x, y$-plane) points in the direction of
greatest increase, it points in the direction in which the graph is
tilted up.  If we project the gradient vector onto the graph, as in
the figure at the top of the next page, it points directly ``uphill''.
Putting it another way, we can \linebreak
	
\newpage

\special{psfile=../ps/922/linear16.ps}
\hfill
\noindent
\begin{minipage}[t]{3.6in}
say that the gradient shows us the ``overall tilt'' of the graph.
There are two parts to this information.  First, the direction of the
gradient tells us which way the graph is tilted.  Second, the length
tells us how steep the graph is.

\quad \ The figure at the left combines all the visual elements we
have introduced to analyze a linear function: contours, graph, and
gradient.  Study it to see how they are related.
 \end{minipage}




\subsection{The Microscope Equation}
\index{microscope equation}

\subsubsection{Local linearity}
\index{local linearity}

Let's return to arbitrary functions of two variables---that is, ones
that are not necessarily linear.\index{function!of two variables}
%
\mar{Local linearity}	
%
First we looked at magnifications of their graphs and contour plots
under a microscope.  We found that the graph becomes a plane, and the
plot becomes a series of parallel, equally-spaced lines.  Next, we saw
that it is precisely the linear functions which have planar graphs and
uniformly parallel contour plots. Hence this function is {\bf locally
  linear}.

Of course, not {\em every\/} function is locally linear, 
%
\mar{Exceptions}	
%
and even a function that {\em is\/} locally linear at most points may
fail to be so at particular points.  We have already seen this with
functions of a single variable in chapter 3.  For example, $g(x) =
x^{2/3}$ is locally linear everywhere {\em except\/} the origin.  It
has a sharp spike\index{spike} there.  The two-variable function $f(x,
y) = (x^2 + y^2)^{1/3}$ has the same sort of spike at the origin.  The
two graphs help make it clear that $g$ is just a slice of~$f$\linebreak
\vspace{-10pt}

 \special{psfile=../ps/922/linear18.ps}
 \special{psfile=../ps/922/linear17.ps}
\vspace{165pt}

\noindent
{\small
\hspace{35pt}
$z = g(x) = x^{2/3}$
\hspace{85pt}
$z = f(x, y) = (x^2 + y^2)^{1/3}$
}

\noindent
(constructed by taking $y = 0$).  (Compare chapter~3.2, pages
\pageref{cusp-begin}--\pageref{cusp-end}.)  The spike is just one
example; there are many other ways that a function can fail to be
locally linear.

Functions that are
%
\mar{The relation between calculus and fractals}	
%
{\em nowhere\/} locally linear are now being used in science to
construct what are called {\bf fractal}\index{fractal} models.
However, calculus does not deal with fractals.  On the contrary, we
remind you of the stipulation first made in chapter 3:
\begin{center}
\begin{picture}(255,40)
%\begin{thicklines}
\put(0,0){\framebox (255,40)
{\shortstack {
{\bf Calculus studies functions that are} \\[2pt]
{\bf locally linear almost everywhere.}
}}}
%\end{thicklines}
\end{picture}
\end{center}

\subsubsection{The microscope equation with two input variables}

If the function $z = f(x, y)$ is locally linear, 
%
\mar{The equation of a microscopic view}	
%
then its graph looks like a plane when we view it under a microscope.
The linear equation that describes that plane is the {\bf microscope
  equation}.\index{microscope equation} Since the plane is part of the
graph of $f$, $f$ itself must determine the form of the microscope
equation.  Let's see how that happens.

The idea is to reduce $f$ to a function of one variable and then use
the microscope equation for one-variable functions (described in
chapter~3.3 and~3.7).  Suppose the microscope is focused at
the point $(x, y) = (a, b)$.  If we fix $y$ (at $y = b$), then $z$
depends on $x$ alone: $z = f(x, b)$.  The microscope equation
%
\mar{\vspace{5pt} The microscope equation in the $x$-direction\ldots}	
%
for this function at $x = a$ is just
	$$
\Delta z \approx \dfrac{\partial f}{\partial x}(a, b) \cdot \Delta x.
	$$
The multiplier\index{multiplier}\index{slope!partial}\index{partial
  slope}\index{partial derivative}\index{derivative!partial}
$\partial f/\partial x$ is the rate of change of $f$ with respect
to~$x$.  We need to write it as a partial derivative because $f$ is a
function of two variables.  Geometrically, the multiplier tells us the
slope of a vertical slice of the graph in the $x$-direction.

Now reverse the roles of $x$ and $y$, fixing $x = a$.  The microscope
equation for the function $z = f(a, y)$
%
\mar{\vspace{13pt} \ldots and in the $y$-direction}	
%
at $y = b$ is
	$$
	\Delta z \approx \dfrac{\partial f}{\partial y}(a, b) \cdot \Delta y.
	$$
The multiplier $\partial f/\partial y$ in this equation is the slope
of a vertical slice of the graph in the $y$-direction.  The slopes of
the two vertical slices are indicated in the microscope window that
appears in the foreground of the following figure.


\newpage

\mbox{}

\vspace{-30pt}
\special{psfile=../ps/921/micro5.ps}
\special{psfile=../ps/921/micro6.ps}
\vspace{268pt}

The separate microscope equations 
%
\mar{From partial changes\\ to total change}	
%
for the $x$- and $y$-directions give us the partial changes in $z$.
However, as we saw when we were defining linear functions
(page~\pageref{altdef}), when we know all the {\em partial\/}
changes,\index{change!partial} we can immediately write down the {\em
  total\/} change:\index{change!total}\index{microscope equation}
\begin{center}
\begin{picture}(240,70) 
	\put(0,-2){\framebox (240,67)
	{\shortstack{
	{\bf The microscope equation}: \\ [6pt]
	$\Delta z \approx 
	\dfrac{\partial f}{\partial x}(a, b) \cdot \Delta x
	+ \dfrac{\partial f}{\partial y}(a, b) \cdot \Delta y$
	}}
}
\end{picture}
\end{center}

As always, the origin of the microscope window corresponds to the
point $(a, b, f(a, b))$ on which the microscope is focused.  The
microscope coordinates $\Delta x$, $\Delta y$, and $\Delta z$ measure
distances from this origin.  For the sake of clarity in the figure
above, we put the origin at one corner of the (three-dimensional)
window.

Incidentally, the function shown above is $f(x, y) = x^3 - 4x - y^2$, 
%
\mar{An example}	
%
and the microscope is focused at $(a, b) = (1.5, -1)$.  Since
	$$
\dfrac{\partial f}{\partial x} = 3 x^2 - 4 = 2.75, \qquad \qquad
		\dfrac{\partial f}{\partial y} = -2 y = 2
	$$
when $(x, y) = (1.5, -1)$, the microscope equation is
	$$
	\Delta z \approx 2.75 \, \Delta x + 2 \, \Delta y.
	$$

%\subsubsection{Linear approximation}
\subsection{Linear Approximation}
\index{linear approximation}
\index{approximation!linear}

The microscope equation 
%
\mar{The microscope equation gives \\a linear approximation}	
%
describes a linear function that approximates the original function
near the point on which the microscope is focused.  It is easy to see
exactly how good the approximation is by comparing contour plots of
the two functions.  This is done below.  In the window on the right,
which shows the highest magnification, the solid contours belong to
the original function $f = x^3 - 4x - y^2$.  They are curved, but only
slightly so.  The dotted contours belong to the linear function
	$$
	(z + 3.625) = 2.75 (x - 1.5) + 2 (y + 1) \quad \mbox{or} \quad
		z = 2.75 \, x + 2 \, y - 5.75.
	$$ 

(This is the microscope equation expressed in terms of the original
                variables $x$, $y$ and $z$ instead of the microscope
                coordinates $\Delta x = x - 1.5$, $\Delta y = y -
                (-1)$ and $\Delta z = z - (-3.625)$.)

%We see here the difference between the {\em defining\/} form of a linear function and the {\em intercept\/} form.)  

The
%
\mar{Comparing the contours \ldots}
%
difference between the two sets of contours shows us just how good the
approximation is.  As you can see, the two functions are almost
indistinguishable near the center of the window, which is the point
$(x, y) = (1.5, -1)$.  As we look farther from the center, we find the
contours of $f$ depart more and more from strict linearity.

\special{psfile=../ps/922/micro7.ps}
%\special{psfile=../ps/921/micro8.ps}

\vspace{158pt}

We can also compare 
%
\mar{\ldots and the graph\\ of the function and its linear approximation} 
%
the {\em graphs\/} of a function and its linear approximation.  The
graph of the linear approximation is a plane, of course.  It is, in
fact, the plane that is tangent to the graph of the function.

In the left figure immediately below we see the tangent plane to
the graph of $z = f(x, y) = x^3 - 4x -y^2$ at the point where $(x, y)
= (1.5, -1)$.  On the right is a magnified view at the point of
tangency.  The graph of $f$ is almost flat.  


\mbox{}

\special{psfile=../ps/923/micro9.ps}
\special{psfile=../ps/923/micro10.ps}
\vspace{132pt}

\noindent
The grid\index{grid line} helps us distinguish it from the tangent
plane, which is solid gray.  Such close agreement between the graph
and the plane demonstrates how good the linear approximation is near
the point.  Notice, however, that the plane diverges from the graph as
we move away from the point of tangency.

At first glance you may not think that the
plane\index{tangent}\index{tangent plane}
%
\mar{Tangent planes}
%
in the figures above is tangent to the surface.  The word ``tangent''
comes from the Latin {\em tangere}, to touch.  We sometimes take this
to mean ``touch at one point'', like the plane in the figure at the
left, below.  More properly, though, two objects are {\bf tangent} if
they have the same direction at a point where they meet.  The plane in
the figures above {\em does\/} meet this condition---as the
microscopic view helps make clear.
 
\special{psfile=../ps/923/micro11.ps}
\special{psfile=../ps/923/micro12.ps}
\vspace{155pt}

There are many different ways 
%
\mar{Elliptic points\ldots}
%
that a tangent plane can intersect a surface.  What happens depends on
the shape of the surface at the point of tangency.  The surface could
bend the same way in all directions at that point, or it could bend up
in some directions and down in others.  In the first case, it will
bend away from its tangent plane, so the two will meet at only one
point.  This is called an {\bf elliptic point},\index{elliptic point}
because the intersection turns into an ellipse if we push the plane in
a bit.  The figure on the right, above, shows what happens.


%\newpage

\mbox{} 


\special{psfile=../ps/923/micro9.ps}
\special{psfile=../ps/923/micro13.ps}
\vspace{135pt}

\enlargethispage*{20pt}
Suppose, on the other hand, that the surface bends up in some
directions but down in others.  Then, in some intermediate directions,
it will not be bending at all.\index{hyperbolic point}
%
\mar{\ldots and\\ hyperbolic points}	
%
In those directions it will meet its tangent plane---which doesn't
bend, either.  Typically, there are two pairs of such directions.  The
surface and the tangent plane then intersect in an~\textsf{X}.  This
always happens at a minimax (or saddle point) on a graph. and it also
happens at the first point we considered on the graph of $z = x^3 - 4x
- y^2$.  This is shown again in the figure on the left above.  It is
called a {\bf hyperbolic point}, because the intersection turns into a
hyperbola if we push the plane a bit---as on the right.  The lines
where the tangent plane itself intersects the surface are called {\bf
  asymptotic lines},\index{asymptotic line} because they are the
asymptotes of the hyperbolas.  (To make it easier to see the
elliptical and hyperbolic intersections we made the surface grid
finer.)

\aside{The curve of intersection between a surface and a shifted
  tangent plane is called the {\em Dupin indicatrix}.\index{Dupin
    indicatrix} The Dupin indicatrix can take many forms besides the
  ones we have described here.  However, at almost all points on
  almost all surfaces it turns out to be an ellipse or a hyperbola.
  More precisely, the indicatrix is {\em approximately\/} an ellipse
  or a hyperbola---in the same way that the surface itself is only
  approximately flat.}

Most points on a surface fall into one of two regions; one region
consists of elliptic points, the other of hyperbolic.
%
\mar{Parabolic points}	
%
Points on the boundary between these two regions are said to be {\bf
  parabolic}.\index{parabolic point} On the graph of $z = x^3 - 4x -
y^2$, if $x < 0$ the point is elliptic; if $x > 0$ it is hyperbolic;
and if $x = 0$, it is parabolic.  Try to confirm this yourself, just
by looking at the surface.

\newpage

\aside{The classification of the points into elliptical, parabolic,
  and hyperbolic types is one of the first steps in studying the
  curvature of a surface.\index{curvature}\index{differential
    geometry}\index{geometry!differential} This is part of {\em
    differential geometry\/}; calculus provides an essential language
  and tool.  Differential geometry is used to model the physical world
  at both the cosmic scale (general
  relativity)\index{relativity!general} and the subatomic (string
  theory).\index{string theory}}


\subsection{The Gradient} 
\index{gradient}

Local linearity is a powerful principle.\index{local linearity}  
%
\mar{Consequences of\\ local linearity}	
%
It says that an arbitrary function looks linear when we view it on a
sufficiently small scale.  In particular, all these statements are
approximately true in a microscope window:
	
$\bullet$ the contours are straight, parallel, and equally spaced;

$\bullet$ the graph is a flat plane (the tangent plane);

$\bullet$ the function has a linear formula (the microscope equation);

$\bullet$ the partial derivatives are the slopes of the graph in the \\
\mbox{} \qquad directions of the axes.

\noindent
There is one more aspect of a linear function 
%
\mar{Extending the gradient to non-linear functions}
%
for us to interpret---the gradient.  Since partial rates become
partial derivatives, we can make the following definition for an
arbitrary locally linear function.\index{$\nabla$}
\begin{center}
\begin{picture}(320,75)		
\put(0,-5) {
  \framebox(320,80){
    \begin{minipage}{290pt}			
      {\bf Definition}.  The {\bf gradient} of a function $z = f(x, y)$ is the vector\index{vector} whose components are the partial derivatives:\index{partial derivative}\index{derivative!partial} 
\linebreak
\vspace{-20pt}
      
	$$
	\mbox{grad $f$} = \nabla f  = (f_x, f_y) = 
		\left( \dfrac{\partial f}{\partial x}, 
		\dfrac{\partial f}{\partial y} \right).
	$$
    \end{minipage}
  }
}
\end{picture}
\end{center}

The partial derivatives are functions, so the gradient varies from
point to point.
%
\mar{The gradient\\ vector field}
%
Thus, gradients form a {\bf vector field}\index{vector field} in the
same sense that a dynamical system\index{dynamical system} does (see
chapter 8).  We draw the gradient vector $(f_x(x, y), f_y(x, y))$ as
an arrow whose tail is at the point $(x, y)$.
\medskip

\noindent
{\bf Example}.  \quad $f(x, y) = x^3 - 4x - y^2$, \quad
	\mar{\vspace{50pt}Contours and gradient vectors together} 
$\mbox{grad $f$} = (3x^2 - 4, -2y)$ 

\special{psfile=../ps/924/grad1.ps}
\special{psfile=../ps/924/grad2.ps}

\newpage

\noindent
There is one thing you should notice about the previous figure: 
%
\mar{The vectors are rescaled for clarity}
%
we drew the gradient vectors much shorter than they actually are.  For
instance, at the origin grad $f$ = $(-4, 0)$, but the arrow {\em as
  drawn\/} is closer to $(-.25, 0)$.  The purpose of
rescaling\index{rescaling} is to keep the vectors out of each other's
way, so the overall pattern easier to see.

The example shows contours as well as gradients so we can see how the
two are related.  The result is very striking.  Even though the
vectors vary in length and direction, and the contours vary in
direction and spacing, the two are related the same way they were for
a {\em linear\/} function (page~\pageref{grad}ff).
%
\mar{The relation between gradients and contours}	
%
First of all, each vector is perpendicular to the level curve that
passes through its tail.  Second, the vectors get longer where the
spacing between level curves gets smaller.  The similarity is no
accident, of course; it is a consequence of local linearity.  We can
summarize and extend our observations in the following theorem.  It is
just a modification of the earlier theorem on the gradient of a linear
function (page~\pageref{gradtheorem}).
\begin{center}
\begin{picture}(360,70)(5.2,0)		
\put(0,0) {
  \framebox(360,70){
    \begin{minipage}{330pt}			
      {\bf Theorem}.\label{gradtheorem2}  The gradient vector field of the function $z = f(x, y)$ is perpendicular to its contour lines.  At each point, the {\em direction\/} of the gradient is the direction in which $z$ increases most rapidly; the {\em length\/} is equal to the maximum rate of increase. 
    \end{minipage}
  }
}
\end{picture}
\end{center}

To see why this theorem is true, 
%
\mar{A proof}	
%
just look in a microscope.  The gradient and the contours become the
gradient and the contours of the linear approximation at the point
where the microscope is focused.  But we already know the theorem is
true for linear functions, so there is nothing more to prove.

There is a direct connection\index{gradient!geometric interpretation} 
%
\mar{The gradient\\ points uphill}
%
between the gradient field of a function $z = f(x, y)$ and its graph.
Since the gradient (which is a vector in the \linebreak
\vspace{-16pt}

\special{psfile=../ps/924/grad3.ps}

\noindent
\begin{minipage}[t]{3.5in}	
$x, y$-plane) points in the direction in which $z$ increases most
  rapidly, it points in the direction in which the graph is tilted up.
  Thus, if we project the gradient onto the graph, as we do in the
  figure at the left, it points directly ``uphill''.  Since $f$ is not
  a linear function, both the steepness and the uphill direction vary
  from point to point.  The gradient vector field also varies; in this
  way, it keeps track of the steepness of the graph and the direction
  of its tilt.
\end{minipage}

\subsection{The Gradient of a Function of Three Variables}

Let's take a brief glance at the gradient of a function $f(x, y,
z)$.\index{function!of three variables}\index{gradient} It has three
components:
	$$
\mbox{grad} f = \nabla f = \left( \dfrac{\partial f}{\partial x},
	\dfrac{\partial f}{\partial y}, \dfrac{\partial f}{\partial z}
		\right) = (f_x, f_y, f_z),
	$$
and it defines a vector field in $x, y, z$-space.  At each point, the
gradient of $f$ is perpendicular to the level set through that point,
and it points in the direction in which $f$ increases most rapidly.

{\bf Example}.  In the two boxes below you can compare the gradient
field of $f(x, y, z) = x^2 + y^2 - z^2$ with its level sets.
%
\mar{Visualizing a three-dimensional vector field}
%
At first glance,\index{vector field} you may not find a clear pattern
to the gradient vectors.  After all, the picture is three-dimensional,
and it is difficult to tell whether an arrow is near the front of the
box or the back.  However, there is a pattern: at the top and the
bottom of the\linebreak

\vspace{-9pt} 

\special{psfile=../ps/925/shells4.ps}
\special{psfile=../ps/924/grad4.ps}
\vspace{210pt}

{\small
\hspace{5pt}
$x^2 + y^2 - z^2 = c$
\hspace{125pt}
$\nabla f = (2x, 2y, -2z)$
}

\vspace{6pt}
\noindent
box, arrows point inward; closer to the middle of the box, they flare
outward.  The lowest values of the function occur along the $z$-axis,
inside the shallow bowls that sit at the top and the bottom of the
box.  The highest values occur outside the ``equatorial belt'' formed
by the outermost level set. This is where the $x, y$-plane meets the
middle of the box.  Notice also that the level sets are symmetric
around the $z$-axis.  The gradient field has the same symmetry, though
this is harder to see.

\subsection{Exercises}


In many of these exercises it will be essential to have a computer
program to make graphs and contour plots of functions of two
variables.
	
\numa Obtain the graph of $z = x^3 - 4x - y^2$ on a domain centered at
the point $(x, y) = (1.5, -1)$, and magnify the graph until it looks
like a plane.

\sub Estimate, by eye, the slopes in the $x$-direction and the
$y$-direction of the plane you found in part (a).  [You should find
  numbers between +2 and +3 in both cases.]

\num Obtain a contour plot of $z = x^3 - 4x - y^2$ on a domain
centered at the point $(x, y) = (1.5, -1)$, and magnify the plot until
the contours look straight, parallel, and equally spaced.  Compare
your results with the plots on page~\pageref{plots1}.

\num Continuation.  The purpose of this exercise is to estimate the
rate of change $\Delta z/\Delta x$ at $(x, y) = (1.5, -1)$, using the
most highly magnified contour plot you constructed in the last
exercise.

\special{psfile=../ps/925/contour1.ps}

\sub What is the horizontal spacing $\Delta x$ between the two
contours closest to the point $(1.5, -1)$?  See the illustration at
the right.

\sub Find the $z$-levels $z_1$ and $z_2$ of those contours, and then
compute $\Delta z = z_2 - z_1$.

\sub Compare the value of $\Delta z/\Delta x$ you now obtain with the
slope in the $x$-direction that you estimated in exercise~1.

\num Continuation.  Repeat all the work of the last exercise, this
time for $\Delta z/\Delta y$.

\subsubsection{Linear functions}
\vspace{-10pt}

\numa Find the $z$-intercept of the graph of the linear function given
by the formula
	$$
	z - 3 = 2 (x - 4) - 3 (y + 1).
	$$
	
\sub  Write the formula for this linear function in intercept form.

\numa Write, in initial-value form, the formula for the linear
function $z = L(x, y)$ for which $\Delta z/\Delta x = 3$, $\Delta
z/\Delta y = -2$, and $L(1, 4) = 0$.

\sub Write the intercept form of the same function.  What is the
$z$-intercept of the graph of $L$?

\numa Suppose $z$ is a linear function of $x$ and $y$, and $\Delta
z/\Delta x = -7$, $\Delta z/\Delta y = 12$.  If $(x, y)$ changes from
$(35, 24)$ to $(33, 33)$, what is the total change in $z$?

\sub Suppose $z = 29$ when $(x, y) = (35, 24)$.  What value does $z$
have when $(x, y) = (33, 33)$?  What value does $z$ have when $(x, y)
= (0, 0)$?

\sub  Write the intercept form of the formula for $z$ in terms of $x$ and $y$.


\num Suppose $z$ is a linear function of $x$ and $y$ for which we have
the following information: \addtolength{\arraycolsep}{8pt}
\renewcommand{\arraystretch}{1.1}
	$$
	\begin{array}{r||r|r|r|r|r|r|r|r}
	x  &  5  &  7  &  0  &     &  4  &  0  &  4  &    \\
	\hline 
	y  &  9  &  1  &  4  &  6  &  7  &  0  &     &  -1 \\
	\hline \hline 
	z  &  2  &     & 12  & -7  &  2  &     &  20 &  20 
	\end{array}
	$$
\addtolength{\arraycolsep}{-8pt}
\renewcommand{\arraystretch}{1}

\sub  Fill in the blanks in this table. 
	
\sub  Write the formula for $z$ in intercept form.

\numa Sketch the graph of $z = x - 2y + 7$ on the domain $0 \leq x
\leq 3$, $0 \leq y \leq 3$.

\sub Determine the slope of this graph in the $x$-direction, and
indicate on your sketch where this slope can be found.

\sub Determine the slope of this graph in the $y$-direction, and
indicate on your sketch where this slope can be found.

\num Continuation.  Draw the gradient vector of $z = x - 2y + 7$ in
the $x, y$-plane, and then lift it up so it sits on the graph you drew
in the previous exercise.  Does the gradient point directly
``uphill''?

\num  Sketch the graph of the linear function $z = L(x, y)$ for which
	$$
\dfrac{\Delta z}{\Delta x} = -1, \qquad \dfrac{\Delta z}{\Delta y} = .6,
		\qquad L(1, 1) = 8.
	$$
Be sure your graph shows clearly the slopes in the $x$-direction and
the $y$-direction, and the $z$-intercept.

\num Continuation.  Draw the gradient of the function $L$ from the
previous exercise, and lift it up so it sits on the graph of $L$ you
drew there.  Does the gradient point directly uphill?

\num What is the equation of the linear function $z = L(x, y)$ whose
graph (a) has a slope of $-4$ in the $x$-direction, (b) a slope of +5
in the $y$-direction, and (c) passes through the point $(x, y, z) =
(2, -9, 0)$?

\num Suppose $z = L(x, y)$ is a linear function whose graph contains
the three points
	$$
	(1, 1, 2), \qquad \quad (0, 5, 4), \qquad \quad (-3, 0, 12).
	$$

\sub Determine the partial rates of change $\Delta z/\Delta x$ and
$\Delta z/\Delta y$.

\sub  Where is the $z$-intercept of the graph?

\sub If $(4, 1, c)$ is a point on the graph, what is the value of $c$?

\sub  Is the point $(2, 2, 4)$ on the graph?  Explain your position.

\special{psfile=../ps/925/linear1.ps}
\special{psfile=../ps/925/linear2.ps}
\vspace{220pt}

\num The figure on the left above is the contour plot of a function $z
= L_1(x, y)$.

\sub What are the values of $L_1(1, 0)$, $L_1(2, 0)$, $L_1(3, 0)$,
$L_1(0, 3)$, $L_1(2, 3)$, and $L_1(0, 0)$?

\sub What are the partial rates $\Delta L_1/\Delta x$ and $\Delta
L_1/\Delta y$?

\num Continuation.  Find the values of $L_1(3, 2)$, $L_1(7, 0)$,
$L_1(7, 7)$, $L_1(1.4, 2.9)$, $L_1(-2, 9)$, $L(-10, -100)$.

\sub Find $x$ so that $L_1(x, 0) = 0$.  Find $y$ so that $L_1(0, y) =
0$.

\num Continuation.  Write the intercept form of the formula for
$L_1(x, y)$.

\numa Find the partial rates $\Delta L_2/\Delta x$ and $\Delta
L_2/\Delta y$ of the linear function $L_2$ whose contour plot is shown
at the right on the previous page.

\sub Obtain the intercept form of the formula for $L_2(x, y)$.

\special{psfile=../ps/925/graph1.ps}
\special{psfile=../ps/925/graph2.ps}
\vspace{190pt}

\num The figures above are the graphs of $L_1$ and $L_2$.  Which is
which?  Explain your choice.  (Note that both graphs are shown from
the same viewpoint.  The $x$-axis is on the right, and the $y$-axis is
in the foreground.  The $z$-axis has no scale on it.)

\num  Determine the gradient vectors of $L_1$ and $L_2$. 

\sub Sketch the gradient vector of each function on the $x,y$-plane
and on its own graph.  Does the gradient point directly uphill in each
case?

\num Suppose the gradient vector of the linear function $p = L(q, r)$
is grad $p = \nabla p = (5, -12)$.  If $L(9, 15) = 17$, what is the
value of $L(11, 11)$?

\numa What is the gradient vector of the function $w = 2 u + 5 v$?

\sub At what point on the circle $u^2 + v^2 = 1$ does $w$ have its
largest value?  What is that value?

\numa Write the formula of a linear function $z = L(x, y)$ whose
gradient vector is grad $z = \nabla z = (-3, 4)$

\sub Using your formula for $L$, calculate the total change in $L$
when $\Delta x = 2$, $\Delta y = 1$.

\numa Continuation; in particular, continue to use your formula for $z
= L(x, y)$.  What is the value of $L(0, 0)$?

\sub What is the maximum value of $z$ on the circle $x^2 + y^2 = 1$?
At what point $(x, y)$ does $z$ achieve that value?

\sub Determine the difference between the maximum and minimum values
of $L$ on the circle $x^2 + y^2 = 1$.

\ans{The difference is 10, independent of the formula you use.}

\numa What value does $z = 7x + 3 y + 31$ have when $x = 5$ and $y =
2$?

\sub  If $x$ increases by 2, 
%
\mar{The concept of \\ a {\em trade-off\/}}
%
how must $y$ change so that the value of $z$ doesn't change.  (The
change in $y$ needed to keep $z$ fixed when $x$ changes is called the
{\bf trade-off}.\index{trade-off} See also chapter~3,
page~\pageref{trade off}.)

\sub  What is the trade-off in $y$ when $x$ increases by $\alpha$?

\sub  What is the trade-off in $x$ when $y$ increases by $\beta$?

\numa Suppose $z = L(x, y)$ is a linear function for which $\Delta
z/\Delta x = 5$ and $\Delta z/\Delta y = -2$.  What is the trade-off
in $y$ when $x$ increases by 50?

\sub  What is the trade-off in $x$ when $y$ increases by 1?

\num Suppose $z = L(x, y)$ is a linear function and suppose the
trade-off in $y$ when $x$ increases by 1 is $-4$.

\sub What is the trade-off in $y$ when $x$ is {\em decreased\/} by 3?

\sub What is the trade-off in $x$ when $y$ is increased by 10?  [Note
  that $x$ and $y$ are reversed here, in comparison to the earlier
  parts of this question.]

\num  Suppose $z = L(x, y)$ is a linear function for which we know
	$$
L(3, 7) = -2, \qquad \quad \dfrac{\Delta z}{\Delta x} = 2.
	$$
Suppose also that the trade-off in $y$ when $x$ increases by 10 is $-4$.

\sub  What is the value of $L(7, 7)$?

\sub  If $L(7, \beta) = -2$, what is the value of $\beta$?

\sub  What is the value of $\Delta z/\Delta y$?

\sub  Write the formula for $L(x, y)$ in intercept form.

\numa Suppose the graph of the linear function $z = L(x, y)$ has a
slope of $-1.5$ in the $x$-direction and $-2.4$ in the $y$-direction.
What is the trade-off between $x$ and $y$?  That is, how much should
$y$ change when $x$ is increased by the amount $\alpha$?

\sub Suppose the partial slopes become $+1.5$ and $+2.4$ in the $x$-
and $y$-directions, respectively.  How does that affect the trade-off?
Explain.

\numa Sketch in the $x, y$-plane the set of points where $z = 2x + 3y
+ 7$ has the value 34.

\sub If $x = 10$ then what value must $y$ have so that the point $(x,
y)$ is on the set in part (a)?

\sub If $x$ increases from 10 to 14, how must $y$ change so that the
point $(x, y)$ stays on the set in part (a)?  In other words, what is
the trade-off?
\medskip

\vspace{-10pt}

\aside{The set in the last question is called a {\em trade-off
    line}.\index{trade-off line} Do you see that it is just a contour
  line by another name?}

\numa  Write, in intercept form, the formula for the linear function
	$$
	w - 4 = 3(x - 2) - 7(y + 1) - 2( z - 5).
	$$

\sub  What is the gradient vector of the linear function in part (a)?

\num Suppose the gradient of the linear function $w = L(x, y, z)$ is
$\nabla w = (1, -1, 4)$.  If $L(3, 0, 5) = 10$, what is the value of
$L(1, 2, 3,)$?

\num Describe the level sets of the function $w = f(x, y, z) = x + y +
z$.


\subsubsection{The microscope equation}
\vspace{-10pt}

\num Find the microscope equation for the function $f(x, y) = 3 x^2 +
4 y^2$ at the point $(x, y) = (2, -1)$.

\numa Continuation.  Use the microscope equation to estimate the
values of $f(1.93, -1.05)$ and $f(2.07, -.99)$
 
\sub Calculate the {\em exact\/} values of the quantities in part (a),
and compare those values with the estimates.  In particular, indicate
how many digits of accuracy the estimates have.

\num Find the microscope equation for the function $f(x, y, z) = x^2 y
\sin{z}$ at the point $(x, y, z) = (1, 1, \pi)$.

\num  Suppose $f(87, 453) = 1254$ and
	$$
\dfrac{\partial f}{\partial x}(87, 453) = -3.4, \qquad \quad
		\dfrac{\partial f}{\partial y}(87, 453) = 4.2.
	$$
Estimate the following values: $f(90, 453)$, $f(87, 450)$, $f(90,
450)$, and $f(100, 500)$.  Explain how you got your estimates.

\numa Continuation.  Find an estimate for $y$ to solve the equation
$f(87, y) = 1250$.

\sub  Find an estimate for $x$ to solve $f(x, 450) = 1275$.

\num Continuation: a trade-off.  Go back to the starting values $x =
87$, $y = 453$, and $f(87, 453) = 1254$.  If $x$ increases from 87 to
88, how should $y$ change to keep the value of the function fixed at
1254?

\numa Suppose $Q(27.3, 31.9) = 15.7$ and $Q(27.9, 31.9) = 15.2$.
Estimate the value of $\partial Q/\partial x(27.6, 31.9)$.

\sub  Estimate the value of $Q(27, 31.9)$.

\num Suppose $S(105, 93) = 10$, $S(110, 93) = 10.7$, $S(105, 95) =
9.3$.  Estimate the value of $S(100, 100)$.  Explain how you made your
estimate.


\num Let $P$ be the point $(x, y, z) =(173, -29, 553)$.  Suppose $f(P)
= 48$ and
	$$
	\dfrac{\partial f}{\partial x}(P) = 7, \qquad \quad
		\dfrac{\partial f}{\partial y}(P) = -2, \qquad \quad
		\dfrac{\partial f}{\partial z}(P) = 5.
	$$
Estimate the value of $f(175, -30, 550)$, and explain what you did.

\numa What is the equation of the tangent plane to the graph of $z =
xy$ at the point $(x, y) = (2, -3)$?

\sub  Which has a higher $z$-intercept: the graph or the tangent plane?

\numa  Suppose the function $H(x, y)$ has the microscope equation
	$$
	\Delta H \approx 2.53 \, \Delta x - 1.19 \, \Delta y
	$$
at the point $(x, y) = (35, 26)$.  Sketch the gradient vector $\nabla\!
H$ at that point.

\sub Pick the point exactly one unit away from $(35, 26)$ at which you
estimate $H$ has the largest possible value.

\ans{At $(35.324, 25.848)$, $H$ is about 2.796 units larger than it is
  at $(35, 26)$.}

\num Write the microscope equation for the function $V(x, y) = x^2 y$
at the point $(x, y) = (25, 10)$.

\numa  Continuation.  
%
\mar{Refer to the discussion of error propagation in chapter~3.4}
%
Suppose a cardboard carton has a square base that is 25 inches on a
side and a height of 10 inches.  If there is an error of $\Delta x$
inches in measuring the base and an error of $\Delta y$ inches in
measuring the height, how much error will there be in the calculated
volume?

\sub  Why is this a continuation of the previous question?

\num Continuation.  Which causes a larger percentage error in the
calculated volume: a 1\% error in the measurement of the length of the
base, or a 1\% error in the measurement of the height?

\vspace{-10pt}

\special{psfile=../ps/925/cone.ps}
\hfill
\begin{minipage}[t]{3.5in}
\num A large basin in the shape of a cone is to be used as a water
reservoir.  If the radius $r$ is 186 meters and the depth $h$ is 31
meters, how much water can the basin hold, in cubic meters?
\end{minipage}

\numa Continuation.  If there were a 3\% error in the measurement of
the radius, how much error would that lead to in the calculation of
the capacity of the basin?

\sub If there were a 5\% error in the measurement of the depth, how
much error would there in the calculated capacity of the basin?

\sub If {\em both\/} errors are present in the measurements, what is
the total error in the calculated capacity of the basin?

\num Continuation.  Suppose the measured radius of the basin ($r =
186$ meters) is assumed to be accurate to within 2\%.  The depth has
been measured at 31 meters.  Is it possible to make that measurement
so accurate that the calculated capacity is known to within 5\%?  How
accurate does the depth measurement have to be?

\num Continuation.  Suppose the accuracy of the radius measurement can
only be guaranteed to be 3\%.  Is it still possible to measure the
depth accurately enough to guarantee that the calculated capacity is
accurate to within 5\%?  Explain.
\vspace{-10pt}

\special{psfile=../ps/925/bowl.ps}

\noindent
\begin{minipage}[t]{3.5in}
\num Let's give the basin the more realistic shape of a parabolic
bowl.  If the radius $r$ is still 186 meters and the depth $h$ is
still 31 meters, what is the capacity of the basin now?  Is it larger
or smaller than the conical basin of the same dimensions?
\end{minipage}

\num Continuation.  Determine the error in the calculated capacity of
the bowl if there were a 3\% error in the measurement of the radius
and, at the same time, a 5\% error in the measurement of the depth.

\num Continuation.  Is it possible for the calculated capacity to be
accurate to within 5\% when the measured radius is accurate to within
2\%?  How accurate does the measurement of the depth have to be to
achieve this?

\num The energy of a certain pendulum whose position is $x$ and
velocity is $v$ can be given by the formula
	$$
	E(x, v) = 1 - \cos{x} + \tfrac{1}{2} v^2.
	$$
Suppose the position of the pendulum is known to be $x = \pi/2$ with a
possible error of 5\% and its velocity is $v = 2$ with possible error
of 10\%.  What is the calculated value of the energy, and how
accurately is that value known?

\numa  A frictionless pendulum conserves energy: 
%
\mar{The conservation\\ of energy leads to\\ a trade-off}
%
as the pendulum moves, the value of $E$ does not change.  Suppose $x =
\pi/2$ and $v = 2$, as in the previous exercise.  When $x$ decreases
by $\pi/180$ (this is 1~degree), does $v$ increase or decrease to
conserve energy?

\sub Approximately how much does $v$ change when $x$ decreases by
$\pi/180$?



\subsubsection{Linear approximations}

\numa Write the linear approximation to $f(x, y) = \sin{x} \cos{y}$ at
the point $(x, y) = (0, \pi/2)$.

\sub Write the equation of the tangent plane to the graph of $f(x, y)$
at the point $(x, y) = (0, \pi/2)$.

\sub  Where does the tangent plane in part (b) meet the $x, y$-plane?

\num  Suppose $w(3, 4) = 2$ while 
	$$
	\dfrac{\partial w}{\partial x}(3, 4) = -1, \qquad \quad
		\dfrac{\partial w}{\partial y} (3, 4) = 3.
	$$  
Write the equation of the tangent plane of $w $ at the point $(3, 4)$.

\numa Write the equation of the tangent plane to the graph of the
function $\varphi(x, y) = 3 x^2 + 7 xy - 2 y^2 - 5x + 3 y$ at an
arbitrary point $(x, y) = (a, b)$.

\sub At what point $(a, b)$ is the tangent plane horizontal?

\sub Magnify the graph of $\varphi$ at the point you found in part (b)
until it looks flat.  Is it also horizontal?
\medskip

The next few exercises concern the Lotka--Volterra differential
equations and their linear approximations at an equilibrium point.
Specifically, consider the {\em bounded growth\/} system from
chapter~4.1 (pages \pageref{L-V bounded
  growth-begin}--\pageref{L-V bounded growth-end}):
\begin{align*}
R' &=  .1\,R  - .00001\,R^2 - .005\,RF,   \\
F' &=  .00004\,RF - .04\,F.
\end{align*}

\vspace{-13pt}

\num Confirm that the system has an equilibrium point at $(R, F) =
(1000, 18)$.  Then obtain the phase portrait of the system near that
point.  What kind of equilibrium is there at $(1000, 18)$?

\num  Obtain the linear approximations of the functions
\begin{align*}
g_1(R, F) &= .1\,R  - .00001\,R^2 - .005\,RF,   \\
g_2(R, F) &= .00004\,RF - .04\,F,
\end{align*}
at the point $(r, F) = (1000, 18)$.  Call them $\ell_1(R, F)$ and
$\ell_2(R, F)$, respectively.

\num  Obtain the phase portrait of the {\bf linear} dynamical system
\begin{align*}
R' &= \ell_1(R, F),   \\
F' &= \ell_2(R, F).
\end{align*}

%\vspace{-8pt}

\sub Does this system have an equilibrium at $(1000, 18)$?  Compare
this phase portrait to the phase portrait of the original {\em
  non\/}-linear system.




\subsubsection{The gradient field}
\vspace{-10pt}

\numa Make a sketch of the gradient vector field of $f(x, y) = x^2 -
y^2$ on the domain $-2 \leq x \leq 2$, $-2 \leq y \leq 2$.

\sub Mark on your sketch where the gradient field indicates the
maximum and the minimum values of $f$ are to be found in the domain.

\num Repeat the previous exercise, using $f(x, y) = 2x + 4 x^2 - x^4 -
y^2$ and the domain $-2 \leq x \leq 2$, $-4 \leq y \leq 4$.  (See
exercises~5 and~6 in the previous section, page~\pageref{quartic}.)
%qqq

\num Continuation.  Add to your sketch in the previous exercise a
contour plot of the function $f(x, y) = 2x + 4 x^2 - x^4 - y^2$ and
confirm that each gradient vector is perpend1 of 82icular to the contour that
passes through its base.  (Note: most vectors have no contour passing
through their bases, so you have to infer the shape and position of
such a contour from the contours that {\em are\/} drawn.)

\num Draw a plausible set of contour lines for the function whose
gradient vector field is plotted on the left below.

\num Draw a plausible gradient vector field for the function whose
contour plot is shown on the right below.  $H$ marks a local maximum
and $L$ a local minimum.

\special{psfile=../ps/925/gradient.ps}
\special{psfile=../ps/925/contour2.ps}


\newpage
%%%%%%%%%%%%%%%%%%%%%%%%%%%%%%%%%%%%%%%%%%%%%%%%%%%%%%%%%%%%%%%%%%%%%%%%%%
%                                                                        %
%                                Section 3                               %
%                                                                        %
%%%%%%%%%%%%%%%%%%%%%%%%%%%%%%%%%%%%%%%%%%%%%%%%%%%%%%%%%%%%%%%%%%%%%%%%%%


\section{Optimization}
\index{optimization}

{\bf Optimization} is the process of making the best choice from a
range of possibilities.  (``Optimum''\index{optimum} is the Latin word
for {\em best}.)
%
\mar{The contexts\\ for optimization}     
%
We are all familiar with optimization in the economic arena: managers
of an enterprise typically seek to to maximize profit or minimize cost
by making conscious choices.  It is perhaps more surprising to learn
that we sometimes use the same language to describe physical
processes.  For instance, the atoms in a molecule are arranged so that
their total energy is minimized.  A light ray travels from one point
to another along the path that takes the least time.  Of course atoms
and photons don't make conscious choices.  Nevertheless, the imagery
of optimization is so vivid and useful that we try to invoke it
whenever we can.

Usually, there is a restriction---called a {\bf constraint}---on the
choices that can be made to achieve to best possible
outcome.\index{constraint}
%
\mar{Constrained optimization}  
%
For instance, consider a factory that makes tennis rackets.  We can
expect that the factory managers are instructed to minimize cost {\em
  while producing a given number of rackets}.  This is their
constraint.  Production cost is a function of many quantities that the
managers can control---the number of workers, the wage scale, and the
cost of the raw materials are just a few.  When the managers choose
values for these quantities that minimize the cost function, they must
be sure those values will also satisfy the constraint.

In mathematical terms, optimization 
%
\mar{Mathematical optimization} 
%
is the process of finding the minimum\index{minimum} or
maximum\index{maximum} value of a function.  The presence of
constraints complicates this task, as you shall see.

\subsection{Visual Inspection}

\special{psfile=../ps/931/OPT1.PS}
\hfill
\begin{minipage}[t]{3.5in}
The maximum value of a function is the highest point on its graph; the
minimum value is the lowest.  Shown at the right is the graph of $z =
x^3 - 4x - y^2$ on the domain

\vspace{-7pt}
\begin{center}
        $-2 \leq x \leq 3 \qquad -2 \leq y \leq 3$.
\end{center}

\vspace{-7pt}

The maximum occurs where $(x, y) = (3, 0)$, and it has the value $z =
15$.  The minimum is $z \approx -12.1$, when $(x, y) \approx (1.2,
3)$.  Confirm this yourself by inspecting the graph and then
calculating $z$.
\end{minipage}

\vspace{4pt}

We'll use the term {\bf extreme}\index{extreme}\index{local
  extreme}\index{extreme!local}
%
\mar{Extremes, and local extremes}      
%
to refer to either a maximum or a minimum.  In the example we see
several {\em local\/} extremes.  These are points where the value of
the function is larger or smaller than it is at any {\em nearby\/}
point.  There are local minima at both left-hand corners of the graph,
at $(-2, -2)$ and $(-2, 3)$.  There is another along the front edge,
near $(1.2, -2)$.  There is a local maximum\index{maximum!local}\index{local maximum} in the interior, near $(-1.2, 0)$.  To
decide which local minimum\index{minimum!local}\index{local minimum}
is the {\em true\/} minimum, we must simply look.  It is the same for
local maxima.

In our example,\index{domain} 
%
\mar{The domain is \\a constraint}      
%
the domain of definition of the function acts as a {\em constraint}.
If we change the domain, the positions of the extreme points can
change.  For example, suppose we change the position of the right-hand
border in stages: first, $x \leq 3$; second, $x \leq 2.5$; third, $x
\leq 2$.  (Since the graph itself doesn't change, we use a grid to
show the part of the graph that satisfies%
\label{optimum1}% 
\linebreak
 
\vspace{12pt}
\special{psfile=../ps/931/OPT2.PS}
\special{psfile=../ps/931/OPT3.PS}
\special{psfile=../ps/931/OPT4.PS}
\vspace{155pt}

\noindent
the constraint in each case below.)  At the start, the maximum is on
the boundary.  When we first move the boundary to the left (from $x =
3$ to $x = 2.5$), 
%
\mar{The maximum moves\\---it even jumps} 
%
the maximum just moves with it.  However, when we move the boundary
farther (from $x = 2.5$ to $x = 2$), the maximum jumps from the
boundary to an interior point\index{interior point} \label{max}(near
$(-1.2, 0)$).  During these changes the minimum is not affected.  It
stays at the same place.

The sudden jump in the maximum 
%
\mar{Catastrophes}
%
is called a {\bf catastrophe}.\index{catastrophe} The figures explain
what happens.  We impose a constraint $x \leq a$, and then we reduce
the value of the parameter~$a$.  At first, the maximum is at the
boundary point $(a, 0)$.  The position of this point changes smoothly
with $a$, causing a gradual drop in the value of the maximum.
Eventually, the maximum reaches the same value as the interior local
maximum.  (This happens when $a = 4/\sqrt{3}$; see the exercises.)  If
$a$ continues to decrease, the local maximum at $(a, 0)$ then has a
lower value than the local maximum at the interior point, so the true
maximum jumps to the interior.

There are many variations on this pattern.  Whenever a function
depends on a parameter the positions of its extremes do, too.  There
are many ways for the position of an extreme to jump suddenly {\em
  while the parameter is changing gradually}.\index{parameter} Any
such jump is called a catastrophe.
        
\aside{Catastrophes make the task of optimization more interesting.
  If the maximum of a certain function gives the optimal solution to a
  problem, and that maximum jumps to a new position, then the optimal
  solution changes radically.  For example, suppose the problem is to
  determine the minimum-energy configuration of atoms in a molecule.
  When the minimum jumps catastrophically, there is a new
  configuration of the same atoms, producing an {\em isomeric\/} form
  of the molecule.
        
  The quest for an optimum does not have to involve mathematical
  tools.  A good example is John Stuart Mill's\index{Mill, J. S.}
  philosophy of ``the greatest good for the greatest number''.  In
  politics a catastrophe is called a revolution.  Even scientific
  research pursues an optimum in raising the question: ``What is the
  best way to explain certain phenomena?''  The consensus in the
  scientific community can change catastrophically, in what is called
  a {\em paradigm shift\/}\index{paradigm shift} or an intellectual
  revolution.  The geological theory of plate tectonics\index{plate
    tectonics} is a familiar example.  Though proposed in the 1920s,
  it was dismissed until the 1960s, when it was suddenly and
  overwhelmingly accepted.}

\vspace{20pt}

\special{psfile=../ps/931/OPT5.PS}
\special{psfile=../ps/931/OPT6.PS}
\special{psfile=../ps/931/OPT7.PS}
\hfill
\begin{minipage}[t]{3.75in}
\subsubsection{Density plots}
We can also solve optimization problems by inspecting density
plots.\index{plot!density}
\label{density1}
Suppose $z$ is the yield from a process that is controlled by two
inputs $x$ and $y$, and
        $$
        z = 3xy - 2 y^2.
        $$
Initially we take $0 \leq x \leq 4$, $0 \leq y \leq 4$.  The maximum
yield is at the darkest spot in the upper density plot.  It occurs on
the right boundary, near $y = 3$.

\quad \ The position of the maximum \index{maximum}is subject to
change if we have to impose further constraints.\index{constraint} For
instance, suppose the resource $y$ is more limited than we first
assumed, requiring us to set $y \leq 2.3$.  With this added constraint
we see that the maximum shifts to the corner $(x, y) = (4, 2.3)$.
(The points shown in a lighter gray are the ones that have been
removed from consideration by the new constraint.)
\end{minipage}

\newpage

\special{psfile=../ps/931/OPT8.PS}
\special{psfile=../ps/931/OPT9.PS}

\noindent
\begin{minipage}[t]{3.6in}
\quad \ Besides limits on individual resources, we are often faced
with a limit on {\em total\/} resources.  In our case, let's suppose
the limit has the form
\[
x + y \leq 4.
\]
That means all points above the line $x + y = 4$ must be removed from
consideration.  (They are shown in lighter gray.)  This new constraint
causes the maximum to shift yet again.  It now appears near the point
$(x, y) = (3, 1)$.
\end{minipage}

\vspace{2pt}

As we add constraints that force the maximum to move, the density at
each new maximum is less than it was at the previous one.  Thus, the
value of the maximum itself decreases.
%
\mar{Constraints\\ reduce optimality}     
%
In other words, each added constraint makes the optimal solution
slightly ``less optimal'' than it had been.  This is only to be
expected.

Of course the extremes may appear in the interior\index{interior point}\index{domain} 
%
\mar{Extremes\\ on the interior\\ of a density plot}        
%
of the domain as well as on the boundary---and density plots can show
this.  Here are the density plots that correspond to the graphs of $z
= x^3 - 4x - y^2$ that we saw on page~\pageref{optimum1}.  The maximum
jumps to the interior\index{catastrophe} when the value of $a$ in the
constraint $x \leq a$\linebreak

\vspace{14pt}
\special{psfile=../ps/931/OPT10.PS}
\special{psfile=../ps/931/OPT11.PS}
\special{psfile=../ps/931/OPT12.PS}
\special{psfile=../ps/931/OPT13.PS}
\special{psfile=../ps/931/OPT14.PS}
\vspace{152pt}

\noindent
drops below $a = 4/\sqrt{3}$.  In all three plots the minimum appears
as the bright spot at the top of the rectangle near where $x = 1$.
(The exact location is $(x, y) = (2/\sqrt{3}, 3)$.)  We can also see a
local minimum at the bottom of the rectangle near $x = 1$.  From the
graph we know there are two more at the left corners of the rectangle,
but they are harder to notice in the density plots.


\subsubsection{Contour plots}

A density plot is useful for showing the general location of the highs
and lows, but we can get a more precise picture by switching to a
contour plot.\index{plot!contour} Let's look at the contour plots of
the same problems we just analyzed using density plots.

\special{psfile=../ps/931/OPT15.PS}
\special{psfile=../ps/931/OPT16.PS}
\special{psfile=../ps/931/OPT17.PS}
\hfill
\begin{minipage}[t]{3.5in}
\quad \ We'll start with the function $z = 3xy - 2 y^2$ we first
considered on page~\pageref{density1} and search for its extremes
subject to the constraints\index{constraint}
\begin{gather*}
0 \leq x,        \\
0 \leq y \leq 2.3, \\
x + y \leq 4, 
\end{gather*}
that we introduced earlier.  From the density plot on
page~\pageref{density1} it was obvious that the values of $z$
\linebreak
\end{minipage}
 
\vspace{-9pt}
\noindent
increase steadily from the upper left corner of the large square to
the upper part of the right side.  In the contour plot, though, we
need some sort of labels to show us where the low and high values of
$z$ are to be found.

\enlargethispage*{5pt}
To locate the extremes within the constrained region, we need to find contours that carry the lowest and the highest values of $z$.  We can do this quite precisely.  
%
\mar{An extreme can occur\\ where a contour\\ is tangent to\\ the boundary}
%
The contour that just passes through the upper left corner---and meets
the region at that point alone---carries the lowest value of $z$.  If
you study the plot you can see that the contour carrying the {\em
  highest\/} value of $z$ also touches the boundary at just a single
point.  It is the contour that is tangent\index{tangent} to the line
$x + y = 4$, and elsewhere lies outside the constrained region.

\special{psfile=../ps/932/opt18.ps}
\special{psfile=../ps/932/opt19.ps}
\special{psfile=../ps/932/opt20.ps}
\hfill
\begin{minipage}[t]{3.5in}
\quad \ Suppose the extreme is in the {\em interior\/}\index{interior
  point} of the constrained region, rather than on the boundary.  For
example, the function $z = x^3 - 4x - y^2$ has an interior maximum
when
\begin{gather*}
-2 \leq x \leq 2        \\
-2 \leq y \leq 3.
\end{gather*}
Around the maximum there is a nest of concentric ovals.  This pattern
is characteristic for an interior extreme.  Notice the local maximum
where a contour is tangent to the boundary line $x = 2$.
\end{minipage}

\vspace{3pt}

These examples demonstrate\index{plot!contour} 
%
\mar{Contour patterns\\ at an extreme}    
%
that there are characteristic patterns that contour lines make near an
extreme point.\index{extreme} One pattern appears along a
constraint\index{constraint} \label{constraint} curve; another pattern
appears at interior points of a domain.  We first 
\linebreak
 
\vspace{12pt}
\special{psfile=../ps/932/opt21.ps}
\special{psfile=../ps/932/opt22.ps}
\vspace{100pt}

\begin{center}
Patterns of contour lines at an extreme point
\end{center}

\vspace{-6pt}
\noindent
met the pattern that appears at an interior\index{interior point}
extreme point when we discussed the `standard' minimum ($x^2 + y^2$)
and maximum ($-x^2 - y^2$) in section~1 (pages \pageref{standard min
  max-begin}--\pageref{standard min max-end}).  {\bf Caution}: the
patterns you see here are ``typical'', but they do not {\em
  guarantee\/} the presence of an extreme point.  The exercises give
you a chance to explore some of the subtleties.

\subsection{Dimension-reducing Constraints}
\index{constraint!dimension-reducing}

Constraints appear frequently in optimization problems.  
%
\mar{How a constraint can reduce the dimension of a problem}    
%
Thus far, however, the constraints we considered have been described
by {\em inequalities}, like $x + y \leq 4$.  Initially, $x$ and $y$
give us the coordinates of a point in a two-dimensional plane.  The
effect of the constraint $x + y \leq 4$ is to restrict the points $(x,
y)$ to just a part of that plane---but it is still a {\em
  two-dimensional\/} part.  Sometimes, though, the constraint is given
by an {\em equality}, like $x + y = 4$.  In that case, the points $(x,
y)$ are restricted to lie on a line---which is a one-dimensional set
in the plane.  The second constraint therefore reduces the dimension
of the problem.

\vspace{18pt}
\special{psfile=../ps/932/constr1.ps}
\special{psfile=../ps/932/constr2.ps}

\newpage

        There is a standard form 
        \mar{A standard form}   
for any constraint that reduces a two-dimensional optimization problem to a one-dimensional problem.  The form is this:
        $$
\mbox{constraint:} \quad g(x, y) = 0.
        $$ 
For instance, the constraint $x + y = 4$ can be written this way by setting $g(x, y) = x + y - 4$.   We'll look at another example in a moment.  

First, though, notice that $g(x, y) = 0$ is one of the contour lines
of the function $g(x, y)$, namely, the contour at level zero.  Since
the contours of $g$ are curves, we often call $g(x, y) = 0$ a {\bf
  constraint curve}.\index{constraint curve}
%
\mar{A comment\\ on dimension}  
%
This curve is {\em one\/}-dimensional, even though it might twist and
turn in a two-dimensional plane.  We say the curve is one-dimensional
because it looks like a straight line under a microscope.  Likewise, a
curved surface is two-dimensional because it looks like a flat plane
under a microscope.\index{dimension!of a curve or surface}

\medskip
\noindent
{\bf Example}.  Find the extreme values of $f(x, y) = x^2 + 8 x y + 3
y^2 - 5 x$ subject to the constraint $g(x, y) = x^2 + y^2 - 4 = 0$.
The constraint curve is a circle of radius~2, and the level curves of
$z = f(x, y)$ form a set of hyperbolas.  We need to find the highest
and the lowest $z$-levels on the constraint curve.

\vspace{7pt}
\special{psfile=../ps/932/constr3.ps}
\vspace{165pt}

The $z$-levels are labelled around the right and the bottom edges of
the figure.
%
\mar{Finding the highest and lowest levels on the constraint curve}
%
It is in the third quadrant, near the point $(-1.4, -1.5)$, that the
constraint curve meets the highest $z$-level. At that point $z$ is
slightly more than 30.  The constraint curve is evidently tangent to
the contour there.  The lowest $z$-level that the constraint curve
meets is about $z = -16$, near the point $(1.6, -1.2)$.  Picture in
your mind how the contours between $z = -10$ and $z = -20$ fit
together.  Can you see that the constraint curve is
tangent\index{tangent} to contour at the minimum?

\newpage

\noindent
{\bf Example, continued}.  
%
\mar{Locate points on the constraint circle with a single variable $t$}
%
There is still more to say about the way the constraint $x^2 + y^2 - 4
= 0$ reduces the dimension of our problem.  First of all, we can
describe the coordinates of any point on the constraint circle by
using the circular functions:
        $$
        x = 2 \cos{t}, \qquad y = 2 \sin{t}.
        $$
We need the factor 2 because the circle has radius 2.  These equations
mean that $x$ and $y$ are now functions of a {\em single\/} variable,
$t$.\index{function!of one variable} Next, consider the function
        $$
        z = f(x, y) = x^2 + 8 x y + 3 y^2 - 5x
        $$
that we seek to optimize {\em subject to the constraint}.  
%
\mar{$z$ becomes a\\ function of $t$}
%
But the constraint makes $x$ and $y$ functions of $t$.  Thus, {\em
  when we take the constraint into account}, $z$ itself becomes a
function of $t$:
\begin{align*}
z &= f(2 \cos{t}, 2 \sin{t})  \\
  &= ( 2 \cos{t} )^2 + 8\, ( 2\cos{t} ) ( 2\sin{t} ) +
                3 ( 2\sin{t} )^2 - 5 ( 2\cos{t} ).
\end{align*} 
The graph of {\em this\/} function is thus just an ordinary curve in
the $t, z$-plane.  It is shown on the right, below.  A value of $t$
determines a point on the circle, as shown in the contour plot on the
left.  (We have taken $t \approx \pi/3$.)  The value of $z$ at that
point then determines the height of the graph on the right.  As you
can see, our chosen value of $t$ puts $z$ near a local maximum.

\vspace{8pt}  
\special{psfile=../ps/932/constr4.ps}
\special{psfile=../ps/932/constr5.ps}

\newpage

\special{psfile=../ps/933/wrap1.ps}
\hfill
\begin{minipage}[t]{3.4in}
{\bf Example, continued}.  Let's look at the graph of $z = x^2 + 8 x y
+ 3 y^2 - 5x$.  The constraint tells us that we should look {\em
  only\/} at the points on the graph that lie above the circle $x^2 +
y^2 - 4 = 0$.  These are the points where the graph intersects the
cylinder\index{cylinder} that you see at the right.  The intersection
is a curve that goes up and down around the cylinder.  At some point
on the curve $z$ has a maximum, and at some other point it has a
minimum.  (In fact, both the maximum and the minimum are visible in
this view.)
\end{minipage} 


\vspace{3pt}

Since the intersection curve lies on the cylinder, 
%
\mar{Unwrap the cylinder}       
%
we can get a better view of the curve if we slit open the cylinder and
unwrap it.  Follow the sequence clockwise from the upper left.  We can
use coordinates to describe the curve on the flattened cylinder.  The
$t$ variable takes us around the cylinder, so it becomes the
horizontal coordinate.  The $z$ variable measures vertical height.
The $z$-range in the figure above is larger than we need: it goes from
$-50$ to $+100$.  In the bottom row on the left we have rescaled the
$z$-axis so it runs from $-20$ to $+35$.  Compare this graph with the
one on the opposite page.
        
\special{psfile=../ps/933/wrap2.ps}
\special{psfile=../ps/933/wrap3.ps}
\special{psfile=../ps/933/wrap5.ps}
\special{psfile=../ps/933/wrap6.ps}
\special{psfile=../ps/933/wrap7.ps}

\newpage

\enlargethispage*{20pt}
The example shows us how a constraint works 
%
\mar{A general view of the dimension-reducing effect of a constraint}   
%
to reduce the dimension of a problem in general.  Suppose we want to
maximize the value of the function $z = f(x, y)$, subject to the
constraint $g(x, y) = 0$.  Then:
\begin{itemize}

\item
{\em Without\/} the constraint, we would look for the highest point on
a {\em two\/}-dimensional surface in three-dimensional space.

\item
{\em With\/} the constraint, we would restrict our search to the
vertical ``curtain''\index{curtain} that lies above the constraint
curve.  The graph intersects this curtain in a curve, so we end up
looking for the highest point on a {\em one\/}-dimensional curve in a
two-dimensional plane.  We can think of this plane as the curtain
after it has been unwrapped and straightened out.

\end{itemize}

\special{psfile=../ps/933/wrap9.ps}
\special{psfile=../ps/933/wrap10.ps}
\special{psfile=../ps/933/wrap11.ps}
\vspace{170pt}

\noindent
This is just a picture of the relation between the function and the
constraint.  We may still have to determine {\em analytically\/} the
form that the function $f(x, y)$ takes when we impose the constraint
$g(x, y) = 0$.  You can find a number of possibilities in the
exercises.


\subsection{Extremes and Critical Points}
\index{extreme}
\index{critical point}

\special{psfile=../ps/933/sample.ps}

\noindent
\begin{minipage}[t]{3.6in}
Suppose that a function $f(x, y)$ has a maximum\index{maximum} or a
minimum\index{minimum} at an {\em interior\/}\index{interior point}
point $(a, b)$.  Suppose also that the function is {\em locally
  linear\/} at $(a, b)$, so we have a ``bowl'' rather than a
``spike''\index{spike} or some other irregularity in the graph.  Then
we must have
        $$
        \mbox{grad} f (a, b) = \left(\dfrac{\partial f}{\partial x}(a, b), 
                \dfrac{\partial f}{\partial y}(a, b) \right)  = (0, 0).
        $$
\end{minipage}

\newpage

\noindent
Here is why.  The gradient\index{gradient} vector $\mbox{grad} f (a, b)$ 
%
\mar{A proof}
%
tells us how the graph of $z = f(x, y)$ is tilted at the point $(a,
b)$.  (We discuss the geometric meaning of the gradient on
page~\pageref{gradtheorem2}.)  At a maximum or a minimum, though, the
graph is {\em not\/} tilted; it must be flat.  Therefore, the gradient
must be the zero vector.

The ideas here carry over to functions 
%
\mar{Moving to\\ an arbitrary number\\ of input variables}  
%
that have any number of input variables.  First, we give a name to a
point where the gradient is zero.
\begin{center}
\begin{picture}(320,60)(5.2,4)          
\put(0,0) {
  \framebox(320,60){
    \begin{minipage}{290pt}                     
      {\bf Definition}.  A {\bf critical point}\index{critical point} of a locally linear function is one where the gradient vector is zero.  Equivalently, all the first partial derivatives of the function are zero.\label{critical}
    \end{minipage}
  }
}
\end{picture}
\end{center}
The observation we just made can now be restated as a theorem that
connects extreme points and critical points.
\begin{center}
\begin{picture}(320,60)(5.2,4)          
\put(0,0) {
  \framebox(320,60){
    \begin{minipage}{290pt}                     
      {\bf Theorem}.  If a locally linear function has a maximum or a minimum at an interior point of its domain, then that point must be a critical point.
    \end{minipage}
  }
}
\end{picture}
\end{center}

The direction of the implication in this theorem is important.  
%
\mar{A statement and\\ its  converse}
%
Here is the theorem, written in a very abbreviated form:
        $$
        \emph{statement}: \qquad 
                \mbox{extreme} \Longrightarrow \mbox{critical}.
        $$
When we reverse the direction of the implication, we get a new
statement, abbreviated the same way:
\[
        \emph{converse}: \qquad \quad
                \mbox{critical} \Longrightarrow \mbox{extreme}.
\]
The converse\index{converse} says that a critical point must be an
extreme point.
%
\mar{The converse of {\em this\/} theorem is not true}
%
But that is just not true.  For example, an ordinary saddle point (a
minimax) is a critical point, but it is not a minimum or a maximum.

The theorem and the observation 
%
\mar{\vspace{9pt}Searching\\ critical points\\ for extremes}    
%
about its converse are both important in the {\em
  optimization\index{optimization} process\/}---that is, the search
for extremes.  Together they offer us the following guidance:

\medskip
$\bullet$ Search for the extremes of a function among its critical points.

\smallskip
$\bullet$ A critical point may be neither a maximum nor a minimum.

To see how we can find extremes by searching among the critical points
of a function, we'll do a few examples.  All the examples use the same
basic idea.  However, as the details get more complicated we bring in
more powerful techniques.
\medskip

\noindent
{\bf Example 1}.  We'll start with the function
        $$
        z = f(x, y) =  x^3 - 4 x - y^2
        $$
we have frequently used as a test case in this chapter.  

The critical points of $f$ are the points that {\em simultaneously\/}
satisfy the two equations\index{derivative!partial}
\begin{align*}
\dfrac{\partial f}{\partial x} &=  3 x^2 - 4 = 0, \\
\dfrac{\partial f}{\partial y} &=  - 2 y = 0.
\end{align*}
        
\special{psfile=../ps/934/contour1.ps}

\noindent
\begin{minipage}[t]{3.5in}
It is clear that $y = 0$ and $x = \pm \sqrt{4/3} \approx \pm 1.1547$.
The two critical points are therefore
        $$
        \left(+\sqrt{4/3}, 0\right) \qquad \mbox{and} \qquad 
                \left(-\sqrt{4/3}, 0\right).
        $$
If you check the contour plot you can see that $(-\sqrt{4/3}, 0)$ is a
local maximum\index{maximum!local} and $(+\sqrt{4/3}, 0)$ is a saddle
point.\index{saddle point}
\end{minipage}

\medskip

\noindent
{\bf Example 2}.  Here is a somewhat more complicated function:
        $$
g(x, y) =  2 x^2 y - y^2 - 4 x^2 + 3 y.
        $$
The critical points are the solutions of the equations
        $$
\dfrac{\partial g}{\partial x} =  4 x y - 8 x = 0, \qquad \qquad
        \dfrac{\partial g}{\partial y} =  2 x^2 - 2 y + 3 = 0.
        $$
Algebraic methods will still work, even though both variables appear in both equations.  For example, we can rewrite ${\partial g/\partial y} = 0$ as
        $$
        2y = 2 x^2 + 3 \qquad \mbox{or} \qquad 
                y = x^2 + \tfrac{3}{2}.
        $$
We can then substitute this expression for $y$ into ${\partial g/\partial x} = 0$ and get
        $$
        4x \left(x^2 + \tfrac{3}{2}\right) - 8x =
        4x \left( x^2 + \tfrac{3}{2} - 2 \right) =
        4x \left( x^2 - \tfrac{1}{2} \right) = 0.
        $$
This implies $x = 0$ or $x = \pm \sqrt{1/2}$.  For each $x$ we can then find the corresponding $y$ by from the equation $y = x^2 + \tfrac{3}{2}$.  

Let's work through this a second time using a geometric approach.  The
equations ${\partial g/\partial x} = 0$ and ${\partial g/\partial y} =
0$ both define curves in the $x, y$-plane.  The curve ${\partial
  g/\partial y} = 0$ is a parabola: $y = x^2 + \tfrac{3}{2}$.  The
equation ${\partial g/\partial x} = 0$ factors as
        $$
        4 x ( y - 2) = 0 \qquad \mbox{or} \qquad x ( y - 2 ) = 0.
        $$

\special{psfile=../ps/933/locus1.ps}
\hfill
\begin{minipage}[t]{3.5in}
Now a product equals 0 precisely when one of its factors equals 0, so
${\partial g/\partial x} = 0$ implies that {\em either\/} $x = 0$ {\em
  or\/} $y - 2 = 0$.  In other words, the ``curve'' ${\partial
  g/\partial x} = 0$ consists of two lines:
\begin{alignat*}{2}
x &= 0,   &\quad &\mbox{a vertical line}, \\
y &= 2,   &\quad &\mbox{a horizontal line}.
\end{alignat*}
The two curves are shown in the figure at the right.  They intersect
in three points.  (The place where the horizontal and vertical lines
cross is {\em not\/} one of the intersection points.)
\end{minipage}


\vspace{3pt}

One of the immediate benefits of the geometric approach is to make it
clear that the $y$-coordinate of a critical point is either 2 or
$\tfrac{3}{2}$.  The critical points are therefore

\special{psfile=../ps/934/contour2.ps}
\hfill
\begin{minipage}[t]{3.5in}
        $$
        \left(-\sqrt{1/2}, 2 \right), \quad
                \left(\vphantom{\sqrt{1/2}}
                        0, \tfrac{3}{2} \right), \quad
                \left(+\sqrt{1/2}, 2 \right).
        $$
A glance at the contour plot of $g$ makes it clear that the first and
third of these are saddle points.  The middle point is a local
maximum.  There are several ways to determine this.  One is to look at
the graph of $g$.  Another is to look at a vertical slice of the graph
through the line $x = 0$.  Then $z = g(0, y) = -y^2 + 3y$.  Here $z$
has a maximum when $y = \tfrac{3}{2}$.
\end{minipage}

\bigskip
        
\aside{We first identified the local maximum of the function $z = x^3
  - 4x - y^2$ on page~\pageref{max}.  At the time, though, we could
  only estimate its position by eye.  Now, however, we can specify its
  location exactly, because we have analytical tools for finding
  critical points.  As the next example shows, these tools are useful
  even when we can't carry out the algebraic manipulations.}

\noindent
{\bf Example 3}.  Here is a function whose critical points we {\em
  can't\/} find using just algebraic computations:
        $$
        z = h(x, y) = x^2 y^2 - x^4 - y^4 - 2x^2 + 5 xy + y
        $$
The equations for the critical points are
\vspace{-10pt}

\special{psfile=../ps/934/contour3.ps}
\special{psfile=../ps/934/contour4.ps}

\noindent
\begin{minipage}[t]{3.8in}
\begin{align*}
\dfrac{\partial h}{\partial x} &=  
                2 x y^2 - 4 x^3 - 4 x + 5 y = 0,  \\[3pt]
\dfrac{\partial h}{\partial y} &=  
                2 x^2 y - 4 y^3 + 5 x + 1 = 0.
\end{align*}
Even though we can't solve the equations algebraically, we can plot
the curves they define by using the contour-plotting program of a
computer.\index{derivative!partial}\index{plot!contour} This is done
at the right.  The curves intersect in three points.  (If you draw the
plots on a large scale, you will find that these are still the only
intersections.)

\quad \ A contour plot of $z = h(x, y)$ itself reveals that the middle
point is a saddle and the outer two are local maxima.  Let's focus on
the maximum in the upper right and determine its position more
precisely.

\quad \ We can always use a microscope.  But if we magnify the contour
plot, we just get a set of nested ovals.  The maximum would lie
somewhere inside the smallest---but we wouldn't know quite where.  By
contrast, the curves ${\partial h/\partial x} = 0$ and ${\partial
  h/\partial y} = 0$ give us a pair of ``crosshairs'' to focus on.
Even with relatively little magnification we can see that the maximum
is at $(1.202\ldots, 1.404\ldots)$.
\end{minipage}

\special{psfile=../ps/934/contour5.ps}
\special{psfile=../ps/934/contour6.ps}
\special{psfile=../ps/934/contour7.ps}

\newpage

\subsection{The Method of Steepest Ascent}
\index{steepest ascent!method of}
\index{method of steepest ascent}

Here is yet another way to find the maximum\index{maximum}  
%
\mar{Walk uphill to get\\ to a maximum}   
%
of a function $z = f(x, y)$.  Imagine that the graph of $f$ is a
landscape that you're standing on.  You want to get to the highest
point. To do that you just walk uphill.  But the uphill direction on
the landscape is given by the gradient vector field of $f$ (see
page~\pageref{gradtheorem2}).\index{gradient!geometric interpretation}
% qqq  
So you move to higher ground by following the gradient field.

In fact, the gradient field defines a 
%
\mar{Trajectories of the gradient dynamical system\ldots}       
%
{\bf dynamical system}\index{dynamical system} of exactly the sort we
studied in chapter~8.  The differential equations are\index{derivative!partial}
\begin{align*}
\dfrac{dx}{dt} &= \dfrac{\partial f}{\partial x}(x, y), \\[3pt]
\dfrac{dy}{dt} &= \dfrac{\partial f}{\partial y}(x, y).
\end{align*}
Because the gradient points uphill, 
%
\mar{\ldots lead to\\ the local maxima}
%
the trajectories\index{trajectory} of this dynamical system also go
uphill.  Trajectories flow to the attractors\index{attractor} of the
system; these are the local maxima of the function.  Furthermore,
since the gradient points in the direction $f$ increases {\em most
  rapidly}, the trajectories follow paths of {\bf steepest ascent} to
the maxima.  This explains the name of the method.
\medskip

{\bf Example 1}.  Let's see how the method of steepest ascent will
find the local maxima of the function
        $$
        z = h(x, y) = x^2 y^2 - x^4 - y^4 - 2x^2 + 5 xy + y
        $$

\special{psfile=../ps/934/ascent1.ps}
\hfill
\begin{minipage}[t]{3.8in}
we considered in the previous example.  The gradient field is
\begin{align*}
\dfrac{dx}{dt} &= 2 x y^2 - 4 x^3 - 4 x + 5 y. \\[3pt]
\dfrac{dy}{dt} &= 2 x^2 y - 4 y^3 + 5 x + 1.
\end{align*}
As you can see at the left, some of the trajectories flow to a local
maximum near $(-1, -1)$, while other flow to the maximum whose
position we determine on the opposite page.  Each attractor has its
own basin of attraction\index{basin of attraction} (as described in
chapter 8).  Therefore, the maximum found by the method of steepest
ascent depends on the initial point of the trajectory.
\end{minipage}


\newpage

If we replace the gradient vectors by their negatives, 
%
\mar{The method of steepest {\em descent}}      
%
the new field will point directly {\em downhill\/}---to the local
minima.  Using the trajectories of the negative gradient field to find
the minima is thus called the method of steepest descent.  In the
following example we use this method to investigate an economic
question.\index{steepest descent!method of}\index{method of steepest
  descent}
\medskip

\special{psfile=../ps/934/distrib1.ps}

\noindent 
{\bf Example 2}.  Manufacturing companies ship their products to
regional warehouses from large distribution centers.  Suppose a
company has regional warehouses at $A$, $B$, and $C$, as shown on the
map at the right.  Where should it put its distribution center $X$ so
as to minimize the total cost of supplying the three regional
warehouses?
\label{distribstart} 

This is a complicated problem that depends on many factors.  For
example, $X$ probably should be put near major roads.  The managers
may also want to choose a location where labor costs are lower.
Certainly, the total distance between the center and the three
warehouses is important.  Let's simply get a {\em first
  approximation\/} to a solution by concentrating on the last factor.
%
\mar{Minimize the\\ total distance}       
%
We will find the position for $X$ that minimizes the total
straight-line distance (the distance ``as the crow flies'') from $X$
to the three points $A$, $B$, and $C$.  The map shows these distances
as three dotted lines.

To describe the various positions we have introduced a coordinate
system in which
\[
A: \ (0, 0), \qquad B: \  (6, 9), \qquad C: \ (10, 2).
\]
The coordinates here are arbitrary.  That is, they don't represent
miles, or kilometers, or any of the usual units of distance---but they
are {\em proportional\/} to the usual units, so we can measure with
them.  If we let the unknown position of $X$ be $(x, y)$, then we seek
to minimize the function
        $$
        S(x, y) = \sqrt{x^2 + y^2} + \sqrt{(x-6)^2 + (y-9)^2} 
                + \sqrt{(x-10)^2 + (y-2)^2}.
        $$

According to the method of steepest descent, we want to find the
attractor of this dynamical system:
\begin{align*}
\dfrac{dx}{dt} &= - \dfrac{x}{\sqrt{x^2 + y^2\vphantom{(y)^2)}}} - 
      \dfrac{x - 6}{\sqrt{(x-6)^2 + (y-9)^2}} -
      \dfrac{x - 10}{\sqrt{(x-10)^2 + (y-2)^2}}, \\
\dfrac{dy}{dt} &= - \dfrac{y}{\sqrt{x^2 + y^2\vphantom{(y)^2)}}} - 
      \dfrac{y - 9}{\sqrt{(x-6)^2 + (y-9)^2}} -
      \dfrac{y - 2}{\sqrt{(x-10)^2 + (y-2)^2}}.
\end{align*}

\newpage

\special{psfile=../ps/934/descent1.ps}
\special{psfile=../ps/934/distrib2.ps}

\mbox{}

\vspace{150pt}

As you can see from the vector field of the dynamical system 
%
\mar{The attractor is\\ a global minimum} 
%
(above left), there is a single attractor, near the point $(6, 4)$.
This implies the total distance function $S(x, y)$ has a single global
minimum---which is what our intuition about the problem would lead us
to expect.
        
To find the position of $X$ more exactly, you can do the
following. First, obtain a solution $(x(t), y(t))$ to the system of
differential equations with an arbitrary initial condition---for
example,
        $$
        x(0)  = 1, \qquad y(0) = 1.
        $$
Then, obtain the coordinates of the 
%
\mar{The attractor is the limit point of a solution} 
%
attractor by evaluating $(x(t), y(t))$ for larger and larger values of
$t$, stopping when the values of $x(t)$ and $y(t)$ stabilize.  You
will find that\index{limit}
        $$
        X = \lim_{t \to \infty}(x(t), y(t)) = (6.22120\ldots, 3.96577\ldots).
        $$  

The important point to note here is that it is not necessary to plot
the vector field---or any other graphic aid, like a contour plot or
graph.
%
\mar{The method needs only a differential equation solver}      
%
You simply need to solve a system of differential equations.  For
example, the values above were found by modifying the computer program
SIRVALUE\index{program!SIRVALUE} we
introduced in chapter~2.  In summary: {\em the method of steepest
  descent (or ascent) requires no graphical tools, but only a basic
  differential equation solver}.\label{distribend}

\newpage

\subsection{Lagrange Multipliers}
\index{Lagrange multiplier}
\index{multiplier!Lagrange}

In searching for the interior\index{interior point} extremes of a
function $f(x, y)$, we have seen that it is helpful to solve the
critical point equations:\index{critical point}
\[
\dfrac{\partial f}{\partial x} = 0, \qquad \quad
                \dfrac{\partial f}{\partial y} = 0.
\]
There is a similar set of equations 
%
\mar{Equations for a constrained extreme} we can use in the search for
the {\em constrained\/}\index{constraint}
extremes\index{extreme}\index{constrained extreme}\index{extreme!constrained} of $f$.  These equations involve a new variable called
a Lagrange multiplier.

So, suppose the function $f(x, y)$ has an extreme on the constraint
curve $g(x, y) = 0$\index{constraint curve} at the point $(a, b)$.
According to the following diagram, the gradient\index{gradient}
vectors $\nabla f$\index{$\nabla$}\index{gradient} and $\nabla g$ must
be parallel at $(a, b)$.
%
\mar{$\nabla f$ and $\nabla g$ must be parallel at a point where $f$ has an extreme along the constraint curve $g = 0$} 
%
(This means they are in the same direction or in opposite directions.)
Here is why.  We already know (see page~\pageref{constraint}) that the
level curve of $f$ that passes through the constrained maximum or
minimum must be tangent to the constraint curve.\index{gradient!geometric interpretation} Now, the gradient vector $\nabla f$ at any
point is perpendicular to the level curve of $f$ through that
point---and the same is true for $g$.  At a point where the level
curves are tangent, the gradients $\nabla f$ and $\nabla g$ are
perpendicular to the {\em same\/} curve, and must therefore be
parallel.
        
        
        
\special{psfile=../ps/934/lagr1.ps}
\vspace{190pt}

Parallel vectors are multiples of each other.  
%
\mar{The multiplier equation}\index{multiplier equation}        
%
Specifically, at a point $(a, b)$ where $\nabla f$ and $\nabla g$ are
parallel, there must be a number $\lambda$\index{$\lambda$} for which
\[
\nabla g (a, b) = \lambda \cdot \nabla f (a, b).
\]
The multiplier $\lambda$ is called a {\bf Lagrange multiplier}.  In
the figure above, $\lambda \approx 1/2$.  If $\nabla g$ and $\nabla f$
were in {\em opposite\/} directions, then $\lambda$ would be negative.

\newpage

\aside{Joseph Louis Lagrange\index{Lagrange, J. L.} (1736--1813) was a
  French mathematician and a younger contemporary of Leonhard
  Euler.\index{Euler, L.}  Like Euler he played an important role in
  making calculus the primary analytical tool for study the physical
  world.  He is particularly noted for his contributions to celestial
  mechanics,\index{celestial mechanics} the field where Isaac
  Newton\index{Newton, I.} first applied the calculus.}

If we write out the multiplier equation using the components of
$\nabla f$ and $\nabla g$, we get\index{derivative!partial}
        $$
\left(\dfrac{\partial g}{\partial x}(a, b), 
                \dfrac{\partial g}{\partial y}(a, b) \right) =
                \lambda \left(\dfrac{\partial f}{\partial x}(a, b), 
                \dfrac{\partial f}{\partial y}(a, b) \right) = 
                \left(\lambda \, \dfrac{\partial f}{\partial x}(a, b), 
                \lambda \, \dfrac{\partial f}{\partial y}(a, b) \right).
        $$
In a vector equation, the vectors are equal component by component.  Thus,
\begin{align*}
\dfrac{\partial g}{\partial x}(a, b) &=
       \lambda \, \dfrac{\partial f}{\partial x}(a, b) \\[4pt]
\dfrac{\partial g}{\partial y}(a, b) &=
       \lambda \, \dfrac{\partial f}{\partial y}(a, b) 
\end{align*}
        
        
Let's return to the main question, 
%
\mar{How can we find a constrained maximum or minimum?} 
%
which can be stated this way: How do we determine where the function
$f(x, y)$ has a maximum or a minimum, subject to the constraint $g(x,
y) = 0$?  If we let $(a, b)$ denote the point we seek, then we see
that $a$ and $b$ satisfy three equations:\index{constraint equation}
\[
\begin{array}{rl}
 \mbox{the constraint equation}: & \hspace{24pt} g(a, b) = 0,     \\[7pt]
 \mbox{the multiplier equations}: &
                \left\{
                \begin{array}{l}
                \dfrac{\partial g}{\partial x}(a, b) = 
                        \lambda \, 
                \dfrac{\partial f}{\partial x}(a, b), \\[14pt]
                \dfrac{\partial g}{\partial y}(a, b) = 
                        \lambda \, 
                        \dfrac{\partial f}{\partial y}(a, b).
                 \end{array}
                 \right.
        \end{array}
\]
In fact, there are {\em three\/} unknowns in these equations: $a$,
$b$, and $\lambda$.  When we solve the three equations for the three
unknowns, we will determine the location of the constrained extreme.
(We'll also have a piece of information we can throw away: the value
of $\lambda$.)
\medskip

\noindent
{\bf Example}.  Find the maximum of $f(x, y) = x^p y^{1-p}$ subject to
the constraint $x + y = c$.  There are two parameters in this problem:
$p$ and $c$.  We assume that $0 < p < 1$ and $0 < c$.  We introduce
parameters to remind you that analytic methods (such as Lagrange
multipliers) are especially valuable in solving problems that depend
on parameters.

We let $g(x, y) = x + y - c$.  Then
       $$
\nabla g = (1, 1), \qquad \quad \nabla f = 
                \left(p x^{p-1} y^{1-p}, (1-p) x^p y^{-p} \right),
        $$
so the three equations we must solve are
\begin{align*}
x + y - c &= 0, \\
        1 &= \lambda p x^{p-1} y^{1-p}, \\
        1 &= \lambda (1-p) x^p y^{-p}.
\end{align*}
Since the second and third equations both equal 1, we can set them
equal to each other:
        $$
        \lambda p x^{p-1} y^{1-p} = \lambda (1-p) x^p y^{-p}.
        $$
We can cancel the two $\lambda$s and combine the powers of $x$ and $y$ to get
        $$
        p x^{-1} y = 1 - p \qquad \mbox{or} \qquad 
                \dfrac{y}{x} = \dfrac{1 - p}{p}.
        $$
According to the first of the three equations, $y = c - x$.  If we substitution this expression in for $y$ into the last equation, we get
        $$
        \dfrac{c - x}{x} = \dfrac{1 - p}{p} \qquad \mbox{or} \qquad 
                p(c - x) = (1 - p) x.
        $$
This equation reduces to $pc = x$, which gives us the $x$-coordinate of the maximum.  To get the $y$-coordinate, we use $y = c - x = c - pc = c(1 - p)$.  To sum up, the maximum is at
        $$
        (x, y) = (cp, c(1-p)) = c(p, 1-p).
        $$
        
\subsection{Exercises}

When searching for an extreme, be sure to zoom in on the graph or plot
you are using as you narrow down the location of the point you seek.

\num Inspect the graph of $z = xy$ to find the maximum value of $z$
subject to the constraints
        $$
x \geq 0, \qquad \quad y \geq 0, \qquad \quad 3 x + 8 y \leq 120.
        $$

\num Inspect the graph of $z = 5 x + 2 y$ to find the minimum value of
$z$ subject to the constraints
        $$
        x \geq 0, \qquad \quad y \geq 0, \qquad \quad xy  \geq 10.
        $$

\num Inspect a contour plot of $z = 3xy - y^2$ to find the maximum
value of $z$ subject to the constraints
        $$
        x \geq 0, \qquad \quad y \geq 0, \qquad \quad x + y \leq 5.
        $$
        
\num Continuation.  Add the constraint $x \leq a$ to the preceding
three, where $a$ is a parameter that takes values between 0 and 5.
Find the maximum value of $z$ subject to all four constraints.
Describe how the position of the constraint depends on the value of
the parameter $a$.

\ans{The maximum is found at $(x, y) = (25/8, 15/8)$, as long as $a
  \geq 25/8$.  When $a < 25/8$, the maximum is at $(x, y) = (a, 5 -
  a)$.}

\num Find the maximum and minimum values of $z = 12 x - 5 y$ when $(x,
y)$ is exactly 1~unit from the origin.

\num  Find the maximum and minimum value 
        \mar{How is the gradient involved here?}
of $z = px + qy$ when $(x, y)$ is exactly 1~unit from the origin.  

\sub At what point is the maximum achieved; at what point is the
minimum achieved?

\num  Use a graph to locate the maximum value of the function 
        $$
        z = 2 x y - 5 x^2 - 7 y^2 + 2 x + 3 y.
        $$
There are no constraints.

\num Use a graph to find the maximum value of $z = 6x + 12 y - x^3 -
y^3$, subject to the constraints
        $$
        x \geq 0, \qquad y \geq 0, \qquad x^2 + y^2 \leq 100.
        $$

\numa Locate the position of the minimum of $x^4 - 2 x^2 - \alpha x +
y^2$ as a function of the parameter $\alpha$.

\sub The position of the minimum jumps catastrophically when $\alpha$
passes through a certain value.  At what value of $\alpha$ does this
happen, and what jump occurs in the minimum?

\numa  Locate the maximum of $x^3 + y^3 - 3x - 3y$ subject to
        $$
        x \leq 3, \qquad \quad y \leq 0, \qquad \quad x + y \leq \beta.
        $$ 
The position of the maximum depends on the value of the parameter
$\beta$, which you can assume lies between 0 and 5.

\sub The position of the maximum jumps catastrophically when $\beta$
crosses a certain threshold value $\beta_0$.  What is $\beta_0$?

\num Find the maximum value of $x^2 y$ in the first quadrant, subject
to the constraint $x + 5 y = 10$.

\num Find the maximum and minimum values of $z = 3x + 4 y$ subject to
the single constraint $x^2 + 4xy + 5y^2 = 10$.

\num  Find all the critical points of the following functions.

\sub  $3 x^2 + 7xy + 2 y^2 + 5x - 6 y + 3$.

\sub $\sin{x} \sin{y}$ \quad on the domain $-4 \leq x \leq 4$, $-4
\leq y \leq 4$.

\sub $\sin{xy}$ \quad on the domain $-4 \leq x \leq 4$, $-4 \leq y
\leq 4$.

\sub  $\exp(x^2 + y^2)$.

\sub  $x^3 + y^3 - 3x - 3 y$.

\sub $x^3 - 3x y^2 - x^2 - y^2$.  [There are four critical points;
  three are saddles.]

\numa  Find the nine critical points of the function
        $$
        C(x, y) = (x^2 + xy + y^2 -1)(x^2 - xy + y^2 - 1).
        $$
Four are minima, four are saddles, and one is a maximum.

\sub Mark the locations of the critical points on a suitable contour
plot of $C(x, y)$.

\num Locate and classify the critical points of the energy integral of
a pendulum:
        $$
        E(x, v) = 1 - \cos{x} + \tfrac{1}{2}v^2.
        $$
        
\sub Compare the {\em critical points\/} of $E$ with the {\em
  equilibrium points\/} of the dynamical system associated with this
energy integral.


\num  Use the method of steepest descent to find the minimum of the function 
        $$
        z = p(x, y) =  e^{2 + y - x^2 - y^2} \sin{x}.
        $$  

\subsubsection{The distribution problem}

The next two exercises are modifications of the distribution problem
on pages \pageref{distribstart}--\pageref{distribend}.  In the example
we assumed that deliveries from the center $X$ to each of the regional
warehouse $A$, $B$, and $C$ happened equally often.  These exercises
assume that deliveries to some warehouses are more frequent than
others.
        
\num Suppose that one truck makes deliveries to $A$ 5 times each week,
while a second truck is used to make 3 deliveries to $B$ and 2 to $C$
each week.  It makes sense to locate $X$ so that the {\em total weekly
  travel\/} is minimized, rather than just the total distance to the
three warehouses.  To get the total weekly travel, we should:

$\bullet$ multiply the distance from $X$ to $A$ by 5;

$\bullet$ multiply the distance from $X$ to $B$ by 3;

$\bullet$ multiply the distance from $X$ to $C$ by 2.\\ 
(Actually, the total {\em round-trip\/} distances are twice these
values, but the proportions would remain the same, so we can use these
numbers.)  Thus, the function to minimize is
        $$
        T(x, y) = 5\sqrt{x^2 + y^2} + 3\sqrt{(x-6)^2 + (y-9)^2} 
                + 2\sqrt{(x-10)^2 + (y-2)^2}.
        $$

\sub  Use the method of steepest descent to find the minimum of $T$.

\sub Compare the location of the distribution center $X$ as determined
by $T$ to its location determined by the function $S$ of example~2 in
the text.  Would you {\em expect\/} the location to change?  In what
direction?  Does the calculated change in position agree with your
intuition?

\num Suppose that deliveries to $A$ are twice as frequent as
deliveries to either $B$ or $C$.  (For example, two trucks make the
round-trip to $A$ each day, but only one truck to $B$ and one to $C$.)
Where should the distribution center $X$ be located in these
circumstances?  Explain how you got your answer.

\ans{$X$ should be at $A$.  Does this surprise you?}

\num A company which has four offices around the country holds an
annual meeting for its top executives.  The location of each office,
and the number of executives at that office, are given in the
following table.  (The coordinates $x$ and $y$ of the position are
given in arbitrary units.)
\begin{center}
\addtolength{\tabcolsep}{6pt}
        \begin{tabular}{ccrr}
        office  & executives    & \multicolumn{1}{c}{$x$} & 
                \multicolumn{1}{c}{$y$} \\
        \hline \\[-10pt]
        $A$     & 32            & 200   & 300   \\
        $B$     & 17            & 1920  & 1100  \\
        $C$     & 20            & 2240  & 450   \\
        $D$     & 41            & 2875  & 1150
        \end{tabular}
\addtolength{\tabcolsep}{-6pt}
\end{center}
Where should the meeting be held if the location depends {\em
  solely\/} on the total travel cost for all the participants?  Assume
that the travel cost, per mile, is the same for every participant.

\subsubsection{The best-fitting line}

Suppose we've taken measurements of two quantities $x$ and $y$, and
obtained the results shown in the table and graph below.  We assume
that $y$ depends on $x$ according to some rule that we don't happen to
know.  In particular, we'd like to know what $y$ is when $x = 5$.  We
have no data.  Can we {\em predict\/} what $y$ should be?

\vspace{15pt}
        
\special{psfile=../ps/934/data1.ps}

\mbox{}
\hspace{55pt}
\addtolength{\tabcolsep}{10pt}
        \begin{tabular}{ll} 
        $x$ & $y$\\ [2pt]
        \hline
        0 & 1 \rule{0pt}{13pt}\\
        1 & 3 \\
        2 & 4 \\
        3 & 3 \\
        4 & 5 \\
        5 & ?
        \end{tabular}
\addtolength{\tabcolsep}{-10pt}
\vspace{20pt}

\special{psfile=../ps/934/data2.ps}

\noindent
\begin{minipage}[t]{3.5in}
\quad \ Here is a common approach to the question.  We assume there is
a simple underlying relation between $y$ and $x$. However, the
measurements that give us the data contain errors or ``noise'' of some
sort that obscure the relationship.  The simplest relation is a linear
function, so we assume that there is a formula $Y = mx + b$ that
describes the connection between $x$ and $y$.
\end{minipage}
\vspace{3pt}


Which line should we choose?  In other words, how should we choose $m$
and $b$?  Since the data points don't lie on a line, there is no
perfect solution.  For any choices, we must expect a difference $e_j$
between the the $j$-th data value $y_j$ and the value $Y_j = m \, j +
b$ predicted by the formula.  These differences are the {\em errors\/}
we assume are present.

A reasonable way to proceed is to {\bf minimize} the total error.
Even this involves choices.  In the figure above, $e_0$ and $e_3$ are
negative, so the ordinary total could be zero, or nearly so, even if
the individual errors were large.  We need a total that {\em
  ignores\/} the signs of the errors.  Here is one:
        $$
        \mbox{absolute error:} \qquad 
                |e_0| + |e_1| + |e_2| + |e_3| + |e_4| = AE.
        $$
Here is another:
        $$
        \mbox{squared error:} \qquad 
                e_0^2 + e_1^2 + e_2^2 + e_3^2 + e_4^2 = SE.
        $$ 
In the following table we compare the data $y_j$ with the calculated values $Y_j$ and the resulting errors $e_j$.
\addtolength{\arraycolsep}{8pt}
        $$
        \begin{array}{cc|cc} 
        x & y & Y & e = y - Y \\ [2pt]
        \hline
        0 & 1 & \rule{0pt}{13pt}b & 1-b    \\
        1 & 3 & \hphantom{1}m+b & 3-\hphantom{1}m-b \\
        2 & 4 & 2m+b & 4-2m-b \\
        3 & 3 & 3m+b & 3-3m-b \\
        4 & 5 & 4m+b & 5-4m-b
        \end{array}
        $$
\addtolength{\arraycolsep}{-8pt}%
To get the values of $AE$ and $SE$, 
%
\mar{The total errors are functions of $m$ and $b$}
%
we take either the absolute values or the squares of the elements of
the rightmost column, and then add.  In particular, the table makes it
clear that both total errors are functions of $m$ and~$b$.  The
absolute error is
        $$
        AE(m, b) = |1-b| + |3-m-b| + |4-2m-b| + |3-3m-b| + |5-4m-b|.
        $$
        
\num Inspect a graph and a contour plot to determine the values of $m$
and $b$ which minimize $AE(m, b)$.

\ans{Remarkable as is may seem, there is an entire line segment of
  solutions to this problem in the $m, b$-plane.  One end of the line
  is near $(m, b) = (.67, 2.3)$, the other is near $(m, b) = (1, 1)$.}

\numa Using a best-fitting line from the previous question, find the
predicted value of $y$ when $x = 5$.

\sub Since there is a range of best-fitting lines, there should be a
range of predicted values for $y$ when $x = 5$.  What is that range?

\num Write down the function $SE(m, b)$ that describes the squared
error in the fit of a straight line to the data given above.

\numa Use a graph and a contour plot to locate the minimum of the
function $SE(m, b)$ from the previous exercise.  Indicate how many
digits of accuracy your answer has.
% m = .8, b = 1.6

\sub Use the method of steepest descent to locate the minimum.  How
many digits of accuracy does {\em this\/} method yield?

\newpage 
%%%%%%%%%%%%%%%%%%%%%%%%%%%%%%%%%%%%%%%%%%%%%%%%%%%%%%%%%%%%%%%%%%%%%%%%%%
%                                                                        %
%                                Section 4                               %
%                                                                        %
%%%%%%%%%%%%%%%%%%%%%%%%%%%%%%%%%%%%%%%%%%%%%%%%%%%%%%%%%%%%%%%%%%%%%%%%%%

\section{Chapter Summary}

\subsection{The Main Ideas}

\begin{itemize}

\item The {\bf graph} of a function of two variables is a
  two-dimensional surface in a three-dimensional space.

\item A function of two variables can also be viewed using a {\bf
  density plot}, a {\bf terraced density plot}, or a {\bf contour
  plot}.  The latter is a set of {\bf level curves} drawn on a flat
  plane.

\item A {\bf contour plot} of a function of three variables is a
  collection of {\bf level surfaces} in three-dimensional space.

\item The graph of a {\bf linear function} is a flat plane, and its
  contour plot consists of straight, parallel, and equally-spaced
  lines.

\item The {\bf gradient} of a linear function is a vector whose
  components are the partial rates of change of the function.

\item Under a {\bf microscope}, the graph of a function of two
  variables becomes a flat plane.  A contour plot turns into a set of
  straight, parallel, and equally-spaced lines.

\item The multipliers in the {\bf microscope equation} for a function
  are its {\bf partial derivatives}:
        $$
        \Delta z = \dfrac{\partial f}{\partial x_1}(a, b)\Delta x_1 + \cdots +
                \dfrac{\partial f}{\partial x_n}(a, b)\Delta x_n.
        $$

\item The {\bf gradient} of a function is a vector whose components
  are the partial derivatives of the function.  Its magnitude and
  direction give the greatest {\bf rate of increase} of the function
  at each point.

\item {\bf Optimization} is a process that involves finding the {\bf
  maximum} or {\bf minimum} value of a function.  There may be {\bf
  constraints} present that limit the scope of the search for an {\bf
  extreme}.

\item Extremes can be found at {\bf critical points}, where all
  partial derivatives of a locally linear function are zero.

\item The {\bf method of steepest ascent} introduces the power of
  dynamical systems into the optimization process.

\end{itemize}

\subsection{Expectations}

\begin{itemize}

\item Using appropriate computer software, you should be able to make
  a \textbf{graph}, a \textbf{terraced density plot}, and a
  \textbf{contour plot} of a function of two variables.

\item Using appropriate graphical representations of a function of two
  variables, you should be able to recognize the \textbf{maxima},
  \textbf{minima}, and \textbf{saddle points}.

\item You should be able to estimate the \textbf{partial rates of
  change} of a function of two variables at a point by zooming in on a
  contour plot.

\item You should be able to recognize the various forms of a
  \textbf{linear function} of several variables and transform the
  representation of the function from one form to another.

\item You should be able to describe the geometric meaning of the
  partial rates of change of a linear function of two variables.

\item You should be able to find the \textbf{gradient} of a function
  of several variables at a point.

\item You should know how the gradient of a function of two variables
  is related to its level curves.


\item You should be able to write the \textbf{microscope equation} for
  a function of two variables at a point.

\item You should be able to use the microscope equation for a function
  of two variables at a point to estimate values of the function at
  nearly points, to find the \textbf{trade-off} in one variable when
  the other changes by a fixed amount, and to estimate errors.

\item You should be able to find the \textbf{linear aprroximation} to
  a function of two variables at a point.

\item You should be able to find the equation of the \textbf{tangent
  plane} to the graph of a function of two variables at a point.

\item You should be able to sketch the \textbf{gradient vector field}
  of a function of two variables in a specified domain.


\item You should be able to sketch a plausible set of contour lines
  for a function whose gradient vecctor field is given; you should be
  able to sketch a plaausible gradient vector field for a ffunction
  whose contour plot is given.

\item You should be able to find the critical points of a function of
  two variables, and you should be able to determine whether a
  critical point is an extreme by inspecting a graph or a contour
  plot.

\item You should be able to find a local maximum of a function of two
  variables by the method of \textbf{steepest descent}.

\item You should be able to find an extreme of a function of two
  variables subject to a constraint either by inspecting a graph or
  contour plot, or by the method of \textbf{Lagrange multipliers}.

\end{itemize}

\cleardoublepage


